
% please use TexLive 2014 or later with the M&C macros freely
% available from tug.org or use any other recent version of LaTeX

\documentclass{book}\usepackage[]{graphicx}\usepackage[]{color}
% maxwidth is the original width if it is less than linewidth
% otherwise use linewidth (to make sure the graphics do not exceed the margin)
\makeatletter
\def\maxwidth{ %
  \ifdim\Gin@nat@width>\linewidth
    \linewidth
  \else
    \Gin@nat@width
  \fi
}
\makeatother

\definecolor{fgcolor}{rgb}{0.345, 0.345, 0.345}
\newcommand{\hlnum}[1]{\textcolor[rgb]{0.686,0.059,0.569}{#1}}%
\newcommand{\hlstr}[1]{\textcolor[rgb]{0.192,0.494,0.8}{#1}}%
\newcommand{\hlcom}[1]{\textcolor[rgb]{0.678,0.584,0.686}{\textit{#1}}}%
\newcommand{\hlopt}[1]{\textcolor[rgb]{0,0,0}{#1}}%
\newcommand{\hlstd}[1]{\textcolor[rgb]{0.345,0.345,0.345}{#1}}%
\newcommand{\hlkwa}[1]{\textcolor[rgb]{0.161,0.373,0.58}{\textbf{#1}}}%
\newcommand{\hlkwb}[1]{\textcolor[rgb]{0.69,0.353,0.396}{#1}}%
\newcommand{\hlkwc}[1]{\textcolor[rgb]{0.333,0.667,0.333}{#1}}%
\newcommand{\hlkwd}[1]{\textcolor[rgb]{0.737,0.353,0.396}{\textbf{#1}}}%
\let\hlipl\hlkwb

\usepackage{framed}
\makeatletter
\newenvironment{kframe}{%
 \def\at@end@of@kframe{}%
 \ifinner\ifhmode%
  \def\at@end@of@kframe{\end{minipage}}%
  \begin{minipage}{\columnwidth}%
 \fi\fi%
 \def\FrameCommand##1{\hskip\@totalleftmargin \hskip-\fboxsep
 \colorbox{shadecolor}{##1}\hskip-\fboxsep
     % There is no \\@totalrightmargin, so:
     \hskip-\linewidth \hskip-\@totalleftmargin \hskip\columnwidth}%
 \MakeFramed {\advance\hsize-\width
   \@totalleftmargin\z@ \linewidth\hsize
   \@setminipage}}%
 {\par\unskip\endMakeFramed%
 \at@end@of@kframe}
\makeatother

\definecolor{shadecolor}{rgb}{.97, .97, .97}
\definecolor{messagecolor}{rgb}{0, 0, 0}
\definecolor{warningcolor}{rgb}{1, 0, 1}
\definecolor{errorcolor}{rgb}{1, 0, 0}
\newenvironment{knitrout}{}{} % an empty environment to be redefined in TeX

\usepackage{alltt}

%the main style; default LibreCaslon font
\usepackage[raggedsec]{morgan2}
\usepackage{morgan-defs}

%to use Times New Roman, instead of LibreCaslon, please uncomment the next line
%\morgansetup{fontsetup=times}

% bibliography
% \usepackage[comma,sort,authoryear]{natbib}         % author-year
\usepackage[comma,sort,numbers]{natbib}            % numbered

% Attempting per-chapter bibliographies. Not working yet.
% \usepackage[sectionbib,compress]{natbib}% provides a bibliography for each chapter
% \usepackage{chapterbib}


% 
% Begin LaTeX packages imported by the authors.
% 

\usepackage{booktabs}    % For awesome table formatting
\usepackage[inline]{enumitem}  % For inline enumerate* lists
\usepackage{gensymb}     % For the \degree command
\usepackage{microtype}   % For (more) beautiful typesetting.
\usepackage{rotating}    % For rotating figures.
\usepackage{wrapfig}     % To wrap text around figures.

% 
% Begin LaTeX macros created by the authors.
% 

\newcommand{\degC}{\degree C}

% 
% Begin R and knitr setup
% 






%
% End author additions
% 


\setcounter{secnumdepth}{2}

\graphicspath{{./figures/}}     % folder for the figures in your book

\PassOptionsToPackage{hyphens}{url}
\usepackage[colorlinks=true,linkcolor=MyDarkBlue,
citecolor=MyDarkBlue,filecolor=MyDarkBlue,urlcolor=MyDarkBlue]{hyperref}

\renewcommand{\UrlBreaks}{\do\.\do\@\do\\\do\/\do\!\do\_\do\|\do\;\do\>\do\]%
\do\)\do\,\do\?\do\&\do\'\do+\do\=\do\#%
\do\a\do\b\do\c\do\d\do\e\do\f\do\g\do\h\do\i\do\j%
\do\k\do\l\do\m\do\n\do\o\do\p\do\q\do\r\do\s\do\t%
\do\u\do\v\do\w\do\x\do\y\do\z\do\A\do\B\do\C\do\D%
\do\E\do\F\do\G\do\H\do\I\do\J\do\K\do\L\do\M\do\N%
\do\O\do\P\do\Q\do\R\do\S\do\T\do\U\do\V\do\W\do\X%
\do\Y\do\Z}
\makeatletter
\g@addto@macro{\UrlBreaks}{\UrlOrds}
\renewcommand{\ALG@name}{\color{black}Algorithm}
\makeatother



\makeindex{}                     % if you are creating an index for your book
\IfFileExists{upquote.sty}{\usepackage{upquote}}{}
\begin{document}
% \SweaveOpts{concordance=TRUE}  % Commented on 17 Oct 2019 to eliminate some errors.

 \frontmatter                   % we'll produce the fm for you; 
    
\def\HALFTITLE{A Framework for Sustainability\\ Thinking: \\%The Multifaceted Nature\\
		%	of Sustainability Challenges:\\
			An Engineering Perspective \\ 
			\textcolor{red}{version $\alpha$ \\}
			\textcolor{black}{\today}}
\def\TITLE{A Framework for Sustainability Thinking \\ %The Multifaceted Nature of Sustainability Challenges\\ 
                        \textcolor{red}{version $\alpha$} \\ \textcolor{black}{\today}}
\def\AUTHORA{Jeremy Van Antwerp}
\def\AFFILIATIONA{Engineering Department, Calvin University}
\def\AUTHORB{Matthew Kuperus Heun}
\def\AFFILIATIONB{Engineering Department, Calvin University}
\def\AUTHORS{\AUTHORA\ and \AUTHORB}
\def\LECTURE{\ \#4}
\def\lcSYNTHESIS{Synthesis Lectures on Sustainable Development}
\def\SYNTHESIS{\MakeUppercase{\textit{\lcSYNTHESIS}}}

\def\EDITOR{xxxx, \textit{yyyy}}


%%%%%%%%%%%% HALF TITLE PAGE 1-2

\thispagestyle{emptyrule}

\halftitle{\HALFTITLE}

\clearpage


%%%%%%%%%%%% FULL TITLE 5

\blankpage

\thispagestyle{emptyrule}
\title{\TITLE}

\vspace*{2pc}
\authorname{\AUTHORA}
\authoraffiliation{\AFFILIATIONA}

\vspace*{1pc}
\authorname{\AUTHORB}
\authoraffiliation{\AFFILIATIONB}

\vfill
\synthesis{\SYNTHESIS\LECTURE}
\morganlogo

\clearpage


%%%%%%%%%%%% ABSTRACT AND KEYWORDS   6

\thispagestyle{emptyrule}

\ABSTRACT
\noindent
%The abstract goes here. 
%The Abstract and the keywords have to fit in this page.
This book is intended to be an introduction to the many challenges of 
sustainability.
It provides basic facts, figures, and information related to sustainability
in a way that is easy to digest.
The focus is on a quantitative understanding of the challenges of sustainability.
The first half of the book develops a framework for sustainability thinking.
The second half considers application areas and personal and corporate responses
to sustainability challenges.
Throughout, the end-of-chapter discussion questions focus on tradeoffs among 
competing goods and the ethical and social implications of decisions related to
sustainability.
This book was written for a university seminar course on sustainability but could
be used in other small-group discussion formats. 
It is intended to be easy to read but hard to digest.
%end ABSTRACT

\keywords{% %% take a look in the index to see what should get added as keywords!
%xxx, yyyy, zz
Sustainability, sustainable development, systems modeling
}

\vfill

\clearpage


%%%%%%%%%%%%% DEDICATION  7-8


%{
%\clearpage
%\thispagestyle{plain}
%
%\vspace*{13pc}\Large\it
%\centerline{To Eric, Jacob, and my parents.}
%
%}
%
%\clearpage
 

%%%%%%%%%%%%%%%%%%%%%%%% TOC %%%%%%%%%%%%%%%%%%

%blankpage

{
\pagestyle{plain}
\tableofcontents
}

\clearpage
 
 
%%%%%%%%%%%%%%%%%%%%%%%% PREFACE %%%%%%%%%%%%%%%%%%%
 
 
\blankpage

{
\chapter*{Preface}
\addcontentsline{toc}{chapter}{\protect\numberline{}{Preface}}
\thispagestyle{plain}
\markboth{PREFACE}{PREFACE}

\noindent
\section*{Who should use this book} 
This book is for anyone who is interested in the many challenges of 
sustainability.
It is an introduction and, as such, does not assume any prior background knowledge,
apart from just a little bit of math.
The book is not intended to be comprehensive but is intended to be easy to read.\\

\section*{How to use this book} 
Read it. Seriously; that's what you do with books.
This book was written to support a one semester-hour university seminar class
on sustainability. The primary focus of class time is in-class discussion, so 
the end-of-chapter discussion questions should provide significant value. 
Therefore, individual chapters, or sets of chapters, can easily be
used to support a wide variety of class structures, types, and levels.
Students should read the text on their own, out of class, to get 
basic information related to  the challenges of sustainability. 
We encourage instructors to use a basic and straightforward quiz on the reading 
through your course management platform to hold students accountable for the 
reading. 
Students should review discussion questions ahead of each class and 
prepare for discussion. \\

\section*{What you should get out of this book} 
This book contains four main contributions.
\begin{itemize}
\item Basic facts, figures, and information related to sustainability. 
We don't assume you start with any prior knowledge but by the time you finish 
reading this book you will know basic facts like carbon dioxide 
concentration in the atmosphere and what are the anthropogenic sources of 
carbon dioxide. 
Furthermore, you should have a \emph{sense of scale} related to sustainability 
information. 
That is, which things are big and important and which are smaller and less relavant. 
To this end, the information is presented \emph{graphically} as much as possible 
to show relative magnitudes.
\item We provide an explanatory framework or scaffold for thinking about sustainability,
primarily in chapters \ref{chap:introduction}-\ref{chap:carbon_intensity}.
\item The end-of-chapter discussion questions get at moral, ethical, philosophical,
and practical aspects of sustainability. 
They don't have %\sout{right} 
objective answers. 
Instead, the questions should help illustrate the role that worldview, values, 
preferences, and \emph{a priori} assumptions have in how we approach sustainability. 
%social pillar?
\item The end-of-chapter project questions range in difficulty from a long
homework problem to a graduate thesis project. 
If graduate study on a sustainability topic is in your future, perhaps you can 
take inspiration from one or more of these project questions. 
Instructors can use one -- or perhaps more -- of these questions for a 
semester-long project, optionally with in-class presentations, for a class.
\end{itemize}
By the end of the book you should see that sustainability is an important, 
urgent, difficult, but solvable problem that requires personal and corporate 
commitments. \\

\section*{Text, organization, and schedule} 
At one chapter a week, the twelve chapters in this book don't quite fill a typical 
university semester. 
Of course, you don't have to be in a university context to use this book. 
A book club or church small group are also appropriate audiences -- sustainability 
is important for everyone. 
However, if you are in a university context, here are a few ways that you could fill the 
remaining weeks of your semester. 
Use class time for project presentations, as mentioned above. 
Perhaps groups of students could each be assigned different project questions and 
at the end of the semester, each could share their results.
Alternatively, a class session or sessions could be devoted to a topic or topics 
your instructor finds personally interesting, like biofuels, grid-scale storage, 
or media coverage of sustainability. 
At Calvin University, our class devotes a week to worldview implications for 
sustainability. 
You could use all or part(s) of \emph{Beyond stewardship: New approaches to 
creation care}~\citep{Warners:2019aa}
or our paper \emph{Current disciplines and worldviews are insufficient
to address sustainability challenges} (2019) as a starting place for discussion. 
% CES website doesn't have 2019 proceedings!
On the other hand, instructors don't need to use all of the book. 
IEEE Xplore makes it easy to assign either a single chapter or a group of chapters 
as supplementary material for many different types of classes.
Lastly, if you read all the way to the end of the Preface, great job! 
Give youself a gold star. 
Now, on to the good stuff...


\vspace*{2pc}
\noindent\AUTHORS\\
\noindent December 2021
}

\clearpage


%%%%%%%%%%%%%%%%%%%%%%%% ACK %%%%%%%%%%%%%%%%%%%


\blankpage

\chapter*{Acknowledgments}
\addcontentsline{toc}{chapter}{\protect\numberline{}{Acknowledgments}}
\thispagestyle{plain}
\markboth{ACKNOWLEDGMENTS}{ACKNOWLEDGMENTS}

\noindent
Thanks to the students of the Sustainability Challenges course and our colleague 
Julie Wildschut for teaching the class in fall 2021.

Thanks to research assistants Larisa Tomeci, Henos Tadesse, and Joshua Broekhuisen
for help with sourcing data and creating figures for the text. 
% was there a 4th student?

Thank you to the many reviewers who provided feedback on early versions to the book:
James VanAntwerp, Julie Wildschut, Becky Haney, James Skillen, Uko Zylstra.

% Jeff Jewett, Shannon Savage, Leonides Murembya, Steve McMullen, Nathan Grawe
% Jennifer VanAntwerp,  Glen VanAntwerp, ...
% Derek Schuurman, Sa BA, Tom Ackerman

Thanks to the Calvin University Board of Trustees and Joel and Linda Zylstra
for funding to create and revise the text.
Thanks also to Tim Koning of Liquid Haulers Maintenance for funding on biofuels research.



\vspace*{2pc}
\noindent\AUTHORS\\
\noindent December 2021
 
\clearpage

\blankpage

                  % we'll need an Abstract and Keywords
								% please see the abs-pref folder for some pdf examples

 \mainmatter
    
% The next command tells RStudio to do "Compile PDF" on book.Rnw,
% instead of this chapter, thereby eliminating the need to switch back to book.Rnw 
% before making the book.
%!TEX root = ../../book.Rnw

\chapter{Introduction}
\label{chap:introduction}

Let’s begin with a story. Once upon a time, a king was challenged to a game of chess by a poor farmer. If he lost, the king would have to pay the farmer a chessboard’s worth of grain, defined as one kernel of grain on the first square, two grains on the second square, four grains on the fourth square, and so on, with each subsequent square receiving double what was on the previous square. The king accepted the wager but, because this is fable, lost the game. A chessboard has eight rows and eight columns, so the king had to pay 264 (less one) grains of rice. Since it takes seven thousand grains of rice to make a pound (~65 mg each), the king owed her 1.195 exograms (one exogram is 1015 kg), which is about 12\% of the weight of everything ever produced by the human species [1]. At US\$6.85/bushel, the grain would be worth about 300 trillion dollars, which is (roughly) ten times more money than exists in the world. 
%Clean up. Maybe kg/person in the world and $/person. 

This story, which goes back at least until the 1200s, illustrates several aspects of the sustainability problem that have been evident for centuries.  Exponential growth is not sustainable. 
Big numbers are hard to grok. % Cognitive biases Decartes’ chiliagon Numerical cognition Subitizing Dunbar number 50-year time frame.
It’s not OK to take giant risks.

The purpose of this book is to summarize the evidence that humans are not living sustainably and to show the characteristics of sustainable societies. The choices and paths to sustainability are bristling with value judgements and moral choices. This book attempts to provide a coherent framework for discussing sustainability that is grounded in a sense of scale. Each chapter is short, with focused discussion questions at the end of each chapter. The questions, which mostly focus on tradeoffs, do not have “right answers,” instead, different answers are indicative of different preferences. Basic data and key concepts start each chapter. This book is not comprehensive; the intent is to equip readers with basic knowledge (informed by scale) so we can grapple with tough moral questions. We intend the book to be easy to read but hard to digest. 

The remainder of this chapter will summarize key themes of the book. Sustainability is existentially important but very difficult for humans to achieve, for several reasons. The numbers and timescales involved are difficult for humans to grasp for a variety of reasons. Sustainability challenges are complex because they involved a variety of interrelated problems from different domain areas. The challenges associated with transitioning to a sustainable existence are complex and interconnected. It is the intention of this book to explore the challenges of sustainability. While we may, at times, point to directions for improved sustainability, it is beyond our scope (and indeed our ability) to provide solutions for all sustainability problems. 

\section{What is sustainability?}
Are humans living "sustainably?" If not, we will -- at some point -- either cease to exist as a species or (at least) experience drastic reductions in our standard of living. Thus, the question of sustainability is of utmost importance. There are many answers to this question. To the novice, many answers may seem to indicate lack of agreement. However, sustainability is almost a self-defining concept. Different answers indicate different assumptions, different priorities, and different boundaries (that is, what is the system under consideration). 

When considering the meaning of sustainability, the two most important questions are "sustaining what?" and "for how long?" The second of these questions is perhaps easier to answer. While humans seem to have difficulty planning for time scales significantly longer than the human lifespan, true sustainability is achieved only if the answer is "indefinitely" or "forever."
The narrowest answer for what needs sustaining is human life and society, which perforce entails those ecosystem services necessary for human health and wellbeing. Beyond these basics, some people view the nonhuman world as having inherent worth or standing and, as such, to also be worth preserving, even if it has no (or negative, for example, smallpox) impact on human continuity. Unfortunately, humans don’t know clearly what pieces of the ecosystem are, in the long run, necessary for survival and which aren’t. For instance, would we be able to survive in a world without dandelions? Maybe. On the other hand, dandelions might be necessary for other organisms we depend on. Environmental science views the ecosystem in its entirety as a web, with all parts depending on all other parts. The ecosystem is viewed then, not as individual components with binary "needed" "not needed" classifications but as a whole that exists on a continuum from "fully functional" to "nonfunctional." The choices and paths to sustainability are bristling with value judgements and moral choices. How much risk do you think is appropriate? What do you value preserving? Different answers to these questions come from different a priori assumptions and values. 

The broadest definitions of sustainability include human-created devices and systems. Humanity could survive without the great pyramids of Giza, although we would have lost great cultural and historical artifacts. On the other hand, we may not survive if we don’t give up coal-fired power plants. Thus, the Three Pillars of sustainability are its environmental, economic, and social aspects. 
Environmental sustainability considers biophysical and thermodynamic constraints and includes issues such as pollution, resource depletion, habitat loss, and biodiversity. Economic sustainability involves questions of profit and loss, wealth management, and macroeconomic policy. Social sustainability comprises human and civil rights, suffering, and personal freedom. This book will discuss issues associated with each of these three areas of sustainability but with a focus on energy and carbon emissions because it is not the only sustainability problem, just the one that needs to be solved most urgently.

The most oft-quoted definition of sustainability comes from the Brundtland report "Sustainable development is development that meets the needs of the present without compromising the ability of future generations to meet their own needs." Note that sustainable development (an oxymoron?!) is not the same as sustainability and “needs” are subjective.

Herman Daly proposed the following three rules for achieving ecological (or environmental) sustainability:
\begin{enumerate}
\item Renewable resources such as fish, soil, and groundwater must be used no faster than the rate at which they regenerate.
\item Nonrenewable resources such as minerals and fossil fuels must be used no faster than renewable substitutes for them can be put into place.
\item Pollution and wastes must be emitted no faster than natural systems can absorb them, recycle them, or render them harmless.
\end{enumerate}

The core idea is that [the rate of] human impact on the ecosystem can be no greater than the ecosystem can absorb, which introduces the idea of the capacity of the ecosystem to provide resources and “take care of” wastes. We will examine the idea of capacity in greater detail in the next chapter.

\section{Impacts and their sources}
Two mathematical expressions, IPAT and Kaya have been used express human impact on the environment and its sources. IPAT expresses “impact” on the environment as the product of human population, affluence, and technology. The Kaya identity is a form of the IPAT equation but limits impact to CO2 emissions. Kaya also helpfully expresses the generic “technology” as the product of energy intensity of the economy and the carbon intensity of energy. 

The first section of this book is organized around a combination of the IPAT and Kaya approach. We keep the general “impact” of IPAT but use the more helpful expression of technology from Kaya. 
\begin{equation}\label{eq:IPARC}
Impact = Population \times Affluence \times Resource intensity \times Contamination intensity
\end{equation}
Impact is a vector quantity (that is, a list) including such things as global warming potential, aquifer depletion, and eutrophication potential. Population is the (scalar) number of people in the world. Affluence is the world per capita GDP. Resource intensity is the vector or list of resources necessary to produce one unit of world GDP. Lastly, impact intensity or contamination intensity is a matrix that lists the impacts of each type of resource.

IPARC is true because it is an identity, and it will provide a useful organizing framework for the following chapters. The terms are equilibrium values, that is they represent static, or steady-state levels, but do not give completely useful information about how changes in any one variable will affect the others. For instance, using resources more efficiently (improving resource intensity), does not lower impact. Instead, it leads to more affluence (see chapter 5). Likewise, it may not be possible to drive the impact(s) per unit of resource to zero because of diminishing returns on efficiency and tradeoffs that exist between different types of impacts.

\section{The mathematics of sustainability}
Returning to the wheat and chessboard problem, for a series of doubling numbers, like the number of grains on a chessboard square: $1, 2, 4, 8, 16, 32, …$, the sum of all the numbers in the series is always one less than the next number in the series. $1+2+4+8 = 15$, which is one less than $16$. $1+2+4+8+16 = 31$, which is one less than $32$. 

Extending the doubling on discrete squares or in discrete time to continuous time is exponential growth. Things that grow exponentially double in a fixed amount of time. An example with a negative growth rate is radioactive decay, where the corresponding time is called the half-life because the amount of time for half of the radioactivity to decay away is constant. When growth rate is constant, the amount of time to double is approximately 70 divided by the growth rate. With a fixed 3\% growth rate, GDP doubles every 23 years. If your investment portfolio has a 7\% return, your assets will double every 10 years. 

In the discrete case, like the wheat and chessboard, the next square is always one more than the sum of all the previous squares. In the continuous case, the amount in the next period (B) is always 1.71 times the sum of all previous amounts (A).
The result for sustainability is that the resources needed by the economy in the next 23 years – for instance, energy – will be more than the total resources (energy) used in all of history. Populations are modeled as having constant growth rates. We actively manage the economy have a constant growth rate (say a target of 3\% GDP growth per year). Because of (exponential) growth in population and affluence, humanity is not on a sustainable path. 

\section{CO2 and Energy}
The approach above (IPARC) is general in that all kinds of impacts can be considered. In this book we will consider many types of impacts and many kinds of sustainability (environmental, economic, and social). However, our focus – in the chapters – is mostly on energy use and CO2 emissions. Reasons for an energy/CO2 focus include the availability of data and that we feel it is (currently) the most urgent sustainability problem. End-of-chapter discussion questions tend to focus more on social aspects of sustainability and worldview issues.


\cleardoublepage
                 % \ref{chap:introduction}
    
% The next command tells RStudio to do "Compile PDF" on book.Rnw,
% instead of this chapter, thereby eliminating the need to switch back to book.Rnw 
% before making the book.
%!TEX root = ../../book.Rnw

\chapter{Planetary Boundaries}
\label{chap:planetary_boundaries}





\cleardoublepage
                     % \ref{chap:planetary_boundaries}
    
% The next command tells RStudio to do "Compile PDF" on book.Rnw,
% instead of this chapter, thereby eliminating the need to switch back to book.Rnw 
% before making the book.
%!TEX root = ../../book.Rnw

\chapter{Population}
\label{chap:population}





\cleardoublepage
                     % \ref{chap:population}
    
% The next command tells RStudio to do "Compile PDF" on book.Rnw,
% instead of this chapter, thereby eliminating the need to switch back to book.Rnw 
% before making the book.
%!TEX root = ../../book.Rnw

\chapter{Affluence}
\label{chap:affluence}





\cleardoublepage
                       % \ref{chap:affluence}
    
% The next command tells RStudio to do "Compile PDF" on book.Rnw,
% instead of this chapter, thereby eliminating the need to switch back to book.Rnw 
% before making the book.
%!TEX root = ../../book.Rnw

\chapter{Energy intensity of the economy}
\label{chap:energy_intensity}





\cleardoublepage
           % \ref{chap:energy_intensity}
    
% The next command tells RStudio to do "Compile PDF" on book.Rnw,
% instead of this chapter, thereby eliminating the need to switch back to book.Rnw 
% before making the book.
%!TEX root = ../../book.Rnw

\chapter{Carbon intensity of energy}
\label{chap:carbon_intensity}





\cleardoublepage
           % \ref{chap:chap:carbon_intensity}
    
% The next command tells RStudio to do "Compile PDF" on book.Rnw,
% instead of this chapter, thereby eliminating the need to switch back to book.Rnw 
% before making the book.
%!TEX root = ../../book.Rnw

\chapter{Households}
\label{chap:households}





\cleardoublepage
                     % \ref{chap:households}
    
% The next command tells RStudio to do "Compile PDF" on book.Rnw,
% instead of this chapter, thereby eliminating the need to switch back to book.Rnw 
% before making the book.
%!TEX root = ../../book.Rnw

\chapter{Transportation}
\label{chap:transportation}





\cleardoublepage
             % \ref{chap:transportation}
    
% The next command tells RStudio to do "Compile PDF" on book.Rnw,
% instead of this chapter, thereby eliminating the need to switch back to book.Rnw 
% before making the book.
%!TEX root = ../../book.Rnw

\chapter{Agriculture}
\label{chap:agriculture}





\cleardoublepage
                   % \ref{chap:agriculture}
    
% The next command tells RStudio to do "Compile PDF" on book.Rnw,
% instead of this chapter, thereby eliminating the need to switch back to book.Rnw 
% before making the book.
%!TEX root = ../../book.Rnw

\chapter{Land use and urban planning}
\label{chap:land_use_urban_planning}





\cleardoublepage
 % \ref{chap:land_use_urban_planning}
    
% The next command tells RStudio to do "Compile PDF" on book.Rnw,
% instead of this chapter, thereby eliminating the need to switch back to book.Rnw 
% before making the book.
%!TEX root = ../../book.Rnw

\chapter{Government}
\label{chap:government}





\cleardoublepage
                     % \ref{chap:government}
    
% The next command tells RStudio to do "Compile PDF" on book.Rnw,
% instead of this chapter, thereby eliminating the need to switch back to book.Rnw 
% before making the book.
%!TEX root = ../../book.Rnw

\chapter{Personal action}
\label{chap:personal_action}





\cleardoublepage
             % \ref{chap:personal_action}

 %   \input{biblio}              % bibliography

    \bibliographystyle{unsrtnat}
    \bibliography{MCBook2021}   %chktex 11
    
    
    \cleardoublepage

 \backmatter                    % back matter
    
%blankpage

\chapter*{Author's Biography}
\markboth{AUTHOR'S BIOGRAPHY}{AUTHOR'S BIOGRAPHY}
\addcontentsline{toc}{chapter}{\protect\numberline{}{Author's Biography}}


\section*{Jeremy Van Antwerp}

% Adjust spacing so the photo looks nice in the paragraph.
\setlength{\intextsep}{-7pt}%
\setlength{\columnsep}{8pt}%
\begin{wrapfigure}{L}{0.25\textwidth}
  \begin{center}
    \includegraphics[width=0.25\textwidth]{headshots/jva-headshot.jpeg}
  \end{center}
\end{wrapfigure}
\textbf{Jeremy Van Antwerp} earned an M.S.\ and Ph.D.\ in chemical engineering 
from the University of Illinois at Urbana-Champaign. 
He is Professor of Engineering (chemical concentration)
at Calvin University in Grand Rapids, MI, USA
and has been a research affiliate and visiting scholar at the Massachusetts 
Institute of Technology, a journal editor, and a consultant in the pharmaceutical 
industry.
His research area is systems and process control.
He holds one US patent and is the author of \emph{Identifcation and control of 
sheet and film processes}~\cite{FeaVB2000}.
His poetry, cartoons, and puzzles have appeared in \emph{IEEE Control Systems}.
He is passionate about excellence and high-impact practices in education.
In 2001 he was awarded a Michigan Traditional Arts apprenticeship grant.


\section*{Matthew Kuperus Heun}

% Adjust spacing so the photo looks nice in the paragraph.
\setlength{\intextsep}{-7pt}%
\setlength{\columnsep}{8pt}%
\begin{wrapfigure}{L}{0.25\textwidth}
  \begin{center}
    \includegraphics[width=0.25\textwidth]{headshots/Heun-pub-photo-headshot}
  \end{center}
\end{wrapfigure}
\textbf{Matthew Kuperus Heun} is Professor of Engineering 
(mechanical concentration)
at Calvin University in Grand Rapids, MI, USA.
He earned an M.S.\ and Ph.D.\ in mechanical engineering from 
the University of Illinois at Urbana-Champaign and 
later worked at NASA's Jet Propulsion Laboratory and at Global Aerospace Corporation. 
He has been a visiting scholar at the Centre for Renewable and Sustainable Energy Studies 
at the University of Stellenbosch, South Africa. 
His long-term research question is 
``What is the relationship between energy and the economy when viewed through the lens of sustainability?''
In addition to scores of articles, he is lead author of 
\emph{Beyond GDP: National accounting in the age of resource depletion}~\citep{Heun:2015aa} 
and a co-editor of
\emph{Beyond Stewardship: New approaches to creation care}~\citep{Warners:2019aa}.
                 % author(s)'s short bio(s)
								% please see the bios folder for some pdf examples
    
    \cleardoublepage

    
% please use TexLive 2014 or later with the M&C macros freely
% available from tug.org or use any other recent version of LaTeX

\documentclass{book}\usepackage[]{graphicx}\usepackage[]{color}
% maxwidth is the original width if it is less than linewidth
% otherwise use linewidth (to make sure the graphics do not exceed the margin)
\makeatletter
\def\maxwidth{ %
  \ifdim\Gin@nat@width>\linewidth
    \linewidth
  \else
    \Gin@nat@width
  \fi
}
\makeatother

\definecolor{fgcolor}{rgb}{0.345, 0.345, 0.345}
\newcommand{\hlnum}[1]{\textcolor[rgb]{0.686,0.059,0.569}{#1}}%
\newcommand{\hlstr}[1]{\textcolor[rgb]{0.192,0.494,0.8}{#1}}%
\newcommand{\hlcom}[1]{\textcolor[rgb]{0.678,0.584,0.686}{\textit{#1}}}%
\newcommand{\hlopt}[1]{\textcolor[rgb]{0,0,0}{#1}}%
\newcommand{\hlstd}[1]{\textcolor[rgb]{0.345,0.345,0.345}{#1}}%
\newcommand{\hlkwa}[1]{\textcolor[rgb]{0.161,0.373,0.58}{\textbf{#1}}}%
\newcommand{\hlkwb}[1]{\textcolor[rgb]{0.69,0.353,0.396}{#1}}%
\newcommand{\hlkwc}[1]{\textcolor[rgb]{0.333,0.667,0.333}{#1}}%
\newcommand{\hlkwd}[1]{\textcolor[rgb]{0.737,0.353,0.396}{\textbf{#1}}}%
\let\hlipl\hlkwb

\usepackage{framed}
\makeatletter
\newenvironment{kframe}{%
 \def\at@end@of@kframe{}%
 \ifinner\ifhmode%
  \def\at@end@of@kframe{\end{minipage}}%
  \begin{minipage}{\columnwidth}%
 \fi\fi%
 \def\FrameCommand##1{\hskip\@totalleftmargin \hskip-\fboxsep
 \colorbox{shadecolor}{##1}\hskip-\fboxsep
     % There is no \\@totalrightmargin, so:
     \hskip-\linewidth \hskip-\@totalleftmargin \hskip\columnwidth}%
 \MakeFramed {\advance\hsize-\width
   \@totalleftmargin\z@ \linewidth\hsize
   \@setminipage}}%
 {\par\unskip\endMakeFramed%
 \at@end@of@kframe}
\makeatother

\definecolor{shadecolor}{rgb}{.97, .97, .97}
\definecolor{messagecolor}{rgb}{0, 0, 0}
\definecolor{warningcolor}{rgb}{1, 0, 1}
\definecolor{errorcolor}{rgb}{1, 0, 0}
\newenvironment{knitrout}{}{} % an empty environment to be redefined in TeX

\usepackage{alltt}

%the main style; default LibreCaslon font
\usepackage[raggedsec]{morgan2}
\usepackage{morgan-defs}

%to use Times New Roman, instead of LibreCaslon, please uncomment the next line
%\morgansetup{fontsetup=times}

% bibliography
% \usepackage[comma,sort,authoryear]{natbib}         % author-year
\usepackage[comma,sort,numbers]{natbib}            % numbered

% Attempting per-chapter bibliographies. Not working yet.
% \usepackage[sectionbib,compress]{natbib}% provides a bibliography for each chapter
% \usepackage{chapterbib}


% 
% Begin LaTeX packages imported by the authors.
% 

\usepackage{booktabs}    % For awesome table formatting
\usepackage[inline]{enumitem}  % For inline enumerate* lists
\usepackage{gensymb}     % For the \degree command
\usepackage{microtype}   % For (more) beautiful typesetting.
\usepackage{rotating}    % For rotating figures.
\usepackage{wrapfig}     % To wrap text around figures.

% 
% Begin LaTeX macros created by the authors.
% 

\newcommand{\degC}{\degree C}

% 
% Begin R and knitr setup
% 






%
% End author additions
% 


\setcounter{secnumdepth}{2}

\graphicspath{{./figures/}}     % folder for the figures in your book

\PassOptionsToPackage{hyphens}{url}
\usepackage[colorlinks=true,linkcolor=MyDarkBlue,
citecolor=MyDarkBlue,filecolor=MyDarkBlue,urlcolor=MyDarkBlue]{hyperref}

\renewcommand{\UrlBreaks}{\do\.\do\@\do\\\do\/\do\!\do\_\do\|\do\;\do\>\do\]%
\do\)\do\,\do\?\do\&\do\'\do+\do\=\do\#%
\do\a\do\b\do\c\do\d\do\e\do\f\do\g\do\h\do\i\do\j%
\do\k\do\l\do\m\do\n\do\o\do\p\do\q\do\r\do\s\do\t%
\do\u\do\v\do\w\do\x\do\y\do\z\do\A\do\B\do\C\do\D%
\do\E\do\F\do\G\do\H\do\I\do\J\do\K\do\L\do\M\do\N%
\do\O\do\P\do\Q\do\R\do\S\do\T\do\U\do\V\do\W\do\X%
\do\Y\do\Z}
\makeatletter
\g@addto@macro{\UrlBreaks}{\UrlOrds}
\renewcommand{\ALG@name}{\color{black}Algorithm}
\makeatother



\makeindex{}                     % if you are creating an index for your book
\IfFileExists{upquote.sty}{\usepackage{upquote}}{}
\begin{document}
% \SweaveOpts{concordance=TRUE}  % Commented on 17 Oct 2019 to eliminate some errors.

 \frontmatter                   % we'll produce the fm for you; 
    
\def\HALFTITLE{A Framework for Sustainability\\ Thinking: \\%The Multifaceted Nature\\
		%	of Sustainability Challenges:\\
			An Engineering Perspective \\ 
			\textcolor{red}{version $\alpha$ \\}
			\textcolor{black}{\today}}
\def\TITLE{A Framework for Sustainability Thinking \\ %The Multifaceted Nature of Sustainability Challenges\\ 
                        \textcolor{red}{version $\alpha$} \\ \textcolor{black}{\today}}
\def\AUTHORA{Jeremy Van Antwerp}
\def\AFFILIATIONA{Engineering Department, Calvin University}
\def\AUTHORB{Matthew Kuperus Heun}
\def\AFFILIATIONB{Engineering Department, Calvin University}
\def\AUTHORS{\AUTHORA\ and \AUTHORB}
\def\LECTURE{\ \#4}
\def\lcSYNTHESIS{Synthesis Lectures on Sustainable Development}
\def\SYNTHESIS{\MakeUppercase{\textit{\lcSYNTHESIS}}}

\def\EDITOR{xxxx, \textit{yyyy}}


%%%%%%%%%%%% HALF TITLE PAGE 1-2

\thispagestyle{emptyrule}

\halftitle{\HALFTITLE}

\clearpage


%%%%%%%%%%%% FULL TITLE 5

\blankpage

\thispagestyle{emptyrule}
\title{\TITLE}

\vspace*{2pc}
\authorname{\AUTHORA}
\authoraffiliation{\AFFILIATIONA}

\vspace*{1pc}
\authorname{\AUTHORB}
\authoraffiliation{\AFFILIATIONB}

\vfill
\synthesis{\SYNTHESIS\LECTURE}
\morganlogo

\clearpage


%%%%%%%%%%%% ABSTRACT AND KEYWORDS   6

\thispagestyle{emptyrule}

\ABSTRACT
\noindent
%The abstract goes here. 
%The Abstract and the keywords have to fit in this page.
This book is intended to be an introduction to the many challenges of 
sustainability.
It provides basic facts, figures, and information related to sustainability
in a way that is easy to digest.
The focus is on a quantitative understanding of the challenges of sustainability.
The first half of the book develops a framework for sustainability thinking.
The second half considers application areas and personal and corporate responses
to sustainability challenges.
Throughout, the end-of-chapter discussion questions focus on tradeoffs among 
competing goods and the ethical and social implications of decisions related to
sustainability.
This book was written for a university seminar course on sustainability but could
be used in other small-group discussion formats. 
It is intended to be easy to read but hard to digest.
%end ABSTRACT

\keywords{% %% take a look in the index to see what should get added as keywords!
%xxx, yyyy, zz
Sustainability, sustainable development, systems modeling
}

\vfill

\clearpage


%%%%%%%%%%%%% DEDICATION  7-8


%{
%\clearpage
%\thispagestyle{plain}
%
%\vspace*{13pc}\Large\it
%\centerline{To Eric, Jacob, and my parents.}
%
%}
%
%\clearpage
 

%%%%%%%%%%%%%%%%%%%%%%%% TOC %%%%%%%%%%%%%%%%%%

%blankpage

{
\pagestyle{plain}
\tableofcontents
}

\clearpage
 
 
%%%%%%%%%%%%%%%%%%%%%%%% PREFACE %%%%%%%%%%%%%%%%%%%
 
 
\blankpage

{
\chapter*{Preface}
\addcontentsline{toc}{chapter}{\protect\numberline{}{Preface}}
\thispagestyle{plain}
\markboth{PREFACE}{PREFACE}

\noindent
\section*{Who should use this book} 
This book is for anyone who is interested in the many challenges of 
sustainability.
It is an introduction and, as such, does not assume any prior background knowledge,
apart from just a little bit of math.
The book is not intended to be comprehensive but is intended to be easy to read.\\

\section*{How to use this book} 
Read it. Seriously; that's what you do with books.
This book was written to support a one semester-hour university seminar class
on sustainability. The primary focus of class time is in-class discussion, so 
the end-of-chapter discussion questions should provide significant value. 
Therefore, individual chapters, or sets of chapters, can easily be
used to support a wide variety of class structures, types, and levels.
Students should read the text on their own, out of class, to get 
basic information related to  the challenges of sustainability. 
We encourage instructors to use a basic and straightforward quiz on the reading 
through your course management platform to hold students accountable for the 
reading. 
Students should review discussion questions ahead of each class and 
prepare for discussion. \\

\section*{What you should get out of this book} 
This book contains four main contributions.
\begin{itemize}
\item Basic facts, figures, and information related to sustainability. 
We don't assume you start with any prior knowledge but by the time you finish 
reading this book you will know basic facts like carbon dioxide 
concentration in the atmosphere and what are the anthropogenic sources of 
carbon dioxide. 
Furthermore, you should have a \emph{sense of scale} related to sustainability 
information. 
That is, which things are big and important and which are smaller and less relavant. 
To this end, the information is presented \emph{graphically} as much as possible 
to show relative magnitudes.
\item We provide an explanatory framework or scaffold for thinking about sustainability,
primarily in chapters \ref{chap:introduction}-\ref{chap:carbon_intensity}.
\item The end-of-chapter discussion questions get at moral, ethical, philosophical,
and practical aspects of sustainability. 
They don't have %\sout{right} 
objective answers. 
Instead, the questions should help illustrate the role that worldview, values, 
preferences, and \emph{a priori} assumptions have in how we approach sustainability. 
%social pillar?
\item The end-of-chapter project questions range in difficulty from a long
homework problem to a graduate thesis project. 
If graduate study on a sustainability topic is in your future, perhaps you can 
take inspiration from one or more of these project questions. 
Instructors can use one -- or perhaps more -- of these questions for a 
semester-long project, optionally with in-class presentations, for a class.
\end{itemize}
By the end of the book you should see that sustainability is an important, 
urgent, difficult, but solvable problem that requires personal and corporate 
commitments. \\

\section*{Text, organization, and schedule} 
At one chapter a week, the twelve chapters in this book don't quite fill a typical 
university semester. 
Of course, you don't have to be in a university context to use this book. 
A book club or church small group are also appropriate audiences -- sustainability 
is important for everyone. 
However, if you are in a university context, here are a few ways that you could fill the 
remaining weeks of your semester. 
Use class time for project presentations, as mentioned above. 
Perhaps groups of students could each be assigned different project questions and 
at the end of the semester, each could share their results.
Alternatively, a class session or sessions could be devoted to a topic or topics 
your instructor finds personally interesting, like biofuels, grid-scale storage, 
or media coverage of sustainability. 
At Calvin University, our class devotes a week to worldview implications for 
sustainability. 
You could use all or part(s) of \emph{Beyond stewardship: New approaches to 
creation care}~\citep{Warners:2019aa}
or our paper \emph{Current disciplines and worldviews are insufficient
to address sustainability challenges} (2019) as a starting place for discussion. 
% CES website doesn't have 2019 proceedings!
On the other hand, instructors don't need to use all of the book. 
IEEE Xplore makes it easy to assign either a single chapter or a group of chapters 
as supplementary material for many different types of classes.
Lastly, if you read all the way to the end of the Preface, great job! 
Give youself a gold star. 
Now, on to the good stuff...


\vspace*{2pc}
\noindent\AUTHORS\\
\noindent December 2021
}

\clearpage


%%%%%%%%%%%%%%%%%%%%%%%% ACK %%%%%%%%%%%%%%%%%%%


\blankpage

\chapter*{Acknowledgments}
\addcontentsline{toc}{chapter}{\protect\numberline{}{Acknowledgments}}
\thispagestyle{plain}
\markboth{ACKNOWLEDGMENTS}{ACKNOWLEDGMENTS}

\noindent
Thanks to the students of the Sustainability Challenges course and our colleague 
Julie Wildschut for teaching the class in fall 2021.

Thanks to research assistants Larisa Tomeci, Henos Tadesse, and Joshua Broekhuisen
for help with sourcing data and creating figures for the text. 
% was there a 4th student?

Thank you to the many reviewers who provided feedback on early versions to the book:
James VanAntwerp, Julie Wildschut, Becky Haney, James Skillen, Uko Zylstra.

% Jeff Jewett, Shannon Savage, Leonides Murembya, Steve McMullen, Nathan Grawe
% Jennifer VanAntwerp,  Glen VanAntwerp, ...
% Derek Schuurman, Sa BA, Tom Ackerman

Thanks to the Calvin University Board of Trustees and Joel and Linda Zylstra
for funding to create and revise the text.
Thanks also to Tim Koning of Liquid Haulers Maintenance for funding on biofuels research.



\vspace*{2pc}
\noindent\AUTHORS\\
\noindent December 2021
 
\clearpage

\blankpage

                  % we'll need an Abstract and Keywords
								% please see the abs-pref folder for some pdf examples

 \mainmatter
    
% The next command tells RStudio to do "Compile PDF" on book.Rnw,
% instead of this chapter, thereby eliminating the need to switch back to book.Rnw 
% before making the book.
%!TEX root = ../../book.Rnw

\chapter{Introduction}
\label{chap:introduction}

Let’s begin with a story. Once upon a time, a king was challenged to a game of chess by a poor farmer. If he lost, the king would have to pay the farmer a chessboard’s worth of grain, defined as one kernel of grain on the first square, two grains on the second square, four grains on the fourth square, and so on, with each subsequent square receiving double what was on the previous square. The king accepted the wager but, because this is fable, lost the game. A chessboard has eight rows and eight columns, so the king had to pay 264 (less one) grains of rice. Since it takes seven thousand grains of rice to make a pound (~65 mg each), the king owed her 1.195 exograms (one exogram is 1015 kg), which is about 12\% of the weight of everything ever produced by the human species [1]. At US\$6.85/bushel, the grain would be worth about 300 trillion dollars, which is (roughly) ten times more money than exists in the world. 
%Clean up. Maybe kg/person in the world and $/person. 

This story, which goes back at least until the 1200s, illustrates several aspects of the sustainability problem that have been evident for centuries.  Exponential growth is not sustainable. 
Big numbers are hard to grok. % Cognitive biases Decartes’ chiliagon Numerical cognition Subitizing Dunbar number 50-year time frame.
It’s not OK to take giant risks.

The purpose of this book is to summarize the evidence that humans are not living sustainably and to show the characteristics of sustainable societies. The choices and paths to sustainability are bristling with value judgements and moral choices. This book attempts to provide a coherent framework for discussing sustainability that is grounded in a sense of scale. Each chapter is short, with focused discussion questions at the end of each chapter. The questions, which mostly focus on tradeoffs, do not have “right answers,” instead, different answers are indicative of different preferences. Basic data and key concepts start each chapter. This book is not comprehensive; the intent is to equip readers with basic knowledge (informed by scale) so we can grapple with tough moral questions. We intend the book to be easy to read but hard to digest. 

The remainder of this chapter will summarize key themes of the book. Sustainability is existentially important but very difficult for humans to achieve, for several reasons. The numbers and timescales involved are difficult for humans to grasp for a variety of reasons. Sustainability challenges are complex because they involved a variety of interrelated problems from different domain areas. The challenges associated with transitioning to a sustainable existence are complex and interconnected. It is the intention of this book to explore the challenges of sustainability. While we may, at times, point to directions for improved sustainability, it is beyond our scope (and indeed our ability) to provide solutions for all sustainability problems. 

\section{What is sustainability?}
Are humans living "sustainably?" If not, we will -- at some point -- either cease to exist as a species or (at least) experience drastic reductions in our standard of living. Thus, the question of sustainability is of utmost importance. There are many answers to this question. To the novice, many answers may seem to indicate lack of agreement. However, sustainability is almost a self-defining concept. Different answers indicate different assumptions, different priorities, and different boundaries (that is, what is the system under consideration). 

When considering the meaning of sustainability, the two most important questions are "sustaining what?" and "for how long?" The second of these questions is perhaps easier to answer. While humans seem to have difficulty planning for time scales significantly longer than the human lifespan, true sustainability is achieved only if the answer is "indefinitely" or "forever."
The narrowest answer for what needs sustaining is human life and society, which perforce entails those ecosystem services necessary for human health and wellbeing. Beyond these basics, some people view the nonhuman world as having inherent worth or standing and, as such, to also be worth preserving, even if it has no (or negative, for example, smallpox) impact on human continuity. Unfortunately, humans don’t know clearly what pieces of the ecosystem are, in the long run, necessary for survival and which aren’t. For instance, would we be able to survive in a world without dandelions? Maybe. On the other hand, dandelions might be necessary for other organisms we depend on. Environmental science views the ecosystem in its entirety as a web, with all parts depending on all other parts. The ecosystem is viewed then, not as individual components with binary "needed" "not needed" classifications but as a whole that exists on a continuum from "fully functional" to "nonfunctional." The choices and paths to sustainability are bristling with value judgements and moral choices. How much risk do you think is appropriate? What do you value preserving? Different answers to these questions come from different a priori assumptions and values. 

The broadest definitions of sustainability include human-created devices and systems. Humanity could survive without the great pyramids of Giza, although we would have lost great cultural and historical artifacts. On the other hand, we may not survive if we don’t give up coal-fired power plants. Thus, the Three Pillars of sustainability are its environmental, economic, and social aspects. 
Environmental sustainability considers biophysical and thermodynamic constraints and includes issues such as pollution, resource depletion, habitat loss, and biodiversity. Economic sustainability involves questions of profit and loss, wealth management, and macroeconomic policy. Social sustainability comprises human and civil rights, suffering, and personal freedom. This book will discuss issues associated with each of these three areas of sustainability but with a focus on energy and carbon emissions because it is not the only sustainability problem, just the one that needs to be solved most urgently.

The most oft-quoted definition of sustainability comes from the Brundtland report "Sustainable development is development that meets the needs of the present without compromising the ability of future generations to meet their own needs." Note that sustainable development (an oxymoron?!) is not the same as sustainability and “needs” are subjective.

Herman Daly proposed the following three rules for achieving ecological (or environmental) sustainability:
\begin{enumerate}
\item Renewable resources such as fish, soil, and groundwater must be used no faster than the rate at which they regenerate.
\item Nonrenewable resources such as minerals and fossil fuels must be used no faster than renewable substitutes for them can be put into place.
\item Pollution and wastes must be emitted no faster than natural systems can absorb them, recycle them, or render them harmless.
\end{enumerate}

The core idea is that [the rate of] human impact on the ecosystem can be no greater than the ecosystem can absorb, which introduces the idea of the capacity of the ecosystem to provide resources and “take care of” wastes. We will examine the idea of capacity in greater detail in the next chapter.

\section{Impacts and their sources}
Two mathematical expressions, IPAT and Kaya have been used express human impact on the environment and its sources. IPAT expresses “impact” on the environment as the product of human population, affluence, and technology. The Kaya identity is a form of the IPAT equation but limits impact to CO2 emissions. Kaya also helpfully expresses the generic “technology” as the product of energy intensity of the economy and the carbon intensity of energy. 

The first section of this book is organized around a combination of the IPAT and Kaya approach. We keep the general “impact” of IPAT but use the more helpful expression of technology from Kaya. 
\begin{equation}\label{eq:IPARC}
Impact = Population \times Affluence \times Resource intensity \times Contamination intensity
\end{equation}
Impact is a vector quantity (that is, a list) including such things as global warming potential, aquifer depletion, and eutrophication potential. Population is the (scalar) number of people in the world. Affluence is the world per capita GDP. Resource intensity is the vector or list of resources necessary to produce one unit of world GDP. Lastly, impact intensity or contamination intensity is a matrix that lists the impacts of each type of resource.

IPARC is true because it is an identity, and it will provide a useful organizing framework for the following chapters. The terms are equilibrium values, that is they represent static, or steady-state levels, but do not give completely useful information about how changes in any one variable will affect the others. For instance, using resources more efficiently (improving resource intensity), does not lower impact. Instead, it leads to more affluence (see chapter 5). Likewise, it may not be possible to drive the impact(s) per unit of resource to zero because of diminishing returns on efficiency and tradeoffs that exist between different types of impacts.

\section{The mathematics of sustainability}
Returning to the wheat and chessboard problem, for a series of doubling numbers, like the number of grains on a chessboard square: $1, 2, 4, 8, 16, 32, …$, the sum of all the numbers in the series is always one less than the next number in the series. $1+2+4+8 = 15$, which is one less than $16$. $1+2+4+8+16 = 31$, which is one less than $32$. 

Extending the doubling on discrete squares or in discrete time to continuous time is exponential growth. Things that grow exponentially double in a fixed amount of time. An example with a negative growth rate is radioactive decay, where the corresponding time is called the half-life because the amount of time for half of the radioactivity to decay away is constant. When growth rate is constant, the amount of time to double is approximately 70 divided by the growth rate. With a fixed 3\% growth rate, GDP doubles every 23 years. If your investment portfolio has a 7\% return, your assets will double every 10 years. 

In the discrete case, like the wheat and chessboard, the next square is always one more than the sum of all the previous squares. In the continuous case, the amount in the next period (B) is always 1.71 times the sum of all previous amounts (A).
The result for sustainability is that the resources needed by the economy in the next 23 years – for instance, energy – will be more than the total resources (energy) used in all of history. Populations are modeled as having constant growth rates. We actively manage the economy have a constant growth rate (say a target of 3\% GDP growth per year). Because of (exponential) growth in population and affluence, humanity is not on a sustainable path. 

\section{CO2 and Energy}
The approach above (IPARC) is general in that all kinds of impacts can be considered. In this book we will consider many types of impacts and many kinds of sustainability (environmental, economic, and social). However, our focus – in the chapters – is mostly on energy use and CO2 emissions. Reasons for an energy/CO2 focus include the availability of data and that we feel it is (currently) the most urgent sustainability problem. End-of-chapter discussion questions tend to focus more on social aspects of sustainability and worldview issues.


\cleardoublepage
                 % \ref{chap:introduction}
    
% The next command tells RStudio to do "Compile PDF" on book.Rnw,
% instead of this chapter, thereby eliminating the need to switch back to book.Rnw 
% before making the book.
%!TEX root = ../../book.Rnw

\chapter{Planetary Boundaries}
\label{chap:planetary_boundaries}





\cleardoublepage
                     % \ref{chap:planetary_boundaries}
    
% The next command tells RStudio to do "Compile PDF" on book.Rnw,
% instead of this chapter, thereby eliminating the need to switch back to book.Rnw 
% before making the book.
%!TEX root = ../../book.Rnw

\chapter{Population}
\label{chap:population}





\cleardoublepage
                     % \ref{chap:population}
    
% The next command tells RStudio to do "Compile PDF" on book.Rnw,
% instead of this chapter, thereby eliminating the need to switch back to book.Rnw 
% before making the book.
%!TEX root = ../../book.Rnw

\chapter{Affluence}
\label{chap:affluence}





\cleardoublepage
                       % \ref{chap:affluence}
    
% The next command tells RStudio to do "Compile PDF" on book.Rnw,
% instead of this chapter, thereby eliminating the need to switch back to book.Rnw 
% before making the book.
%!TEX root = ../../book.Rnw

\chapter{Energy intensity of the economy}
\label{chap:energy_intensity}





\cleardoublepage
           % \ref{chap:energy_intensity}
    
% The next command tells RStudio to do "Compile PDF" on book.Rnw,
% instead of this chapter, thereby eliminating the need to switch back to book.Rnw 
% before making the book.
%!TEX root = ../../book.Rnw

\chapter{Carbon intensity of energy}
\label{chap:carbon_intensity}





\cleardoublepage
           % \ref{chap:chap:carbon_intensity}
    
% The next command tells RStudio to do "Compile PDF" on book.Rnw,
% instead of this chapter, thereby eliminating the need to switch back to book.Rnw 
% before making the book.
%!TEX root = ../../book.Rnw

\chapter{Households}
\label{chap:households}





\cleardoublepage
                     % \ref{chap:households}
    
% The next command tells RStudio to do "Compile PDF" on book.Rnw,
% instead of this chapter, thereby eliminating the need to switch back to book.Rnw 
% before making the book.
%!TEX root = ../../book.Rnw

\chapter{Transportation}
\label{chap:transportation}





\cleardoublepage
             % \ref{chap:transportation}
    
% The next command tells RStudio to do "Compile PDF" on book.Rnw,
% instead of this chapter, thereby eliminating the need to switch back to book.Rnw 
% before making the book.
%!TEX root = ../../book.Rnw

\chapter{Agriculture}
\label{chap:agriculture}





\cleardoublepage
                   % \ref{chap:agriculture}
    
% The next command tells RStudio to do "Compile PDF" on book.Rnw,
% instead of this chapter, thereby eliminating the need to switch back to book.Rnw 
% before making the book.
%!TEX root = ../../book.Rnw

\chapter{Land use and urban planning}
\label{chap:land_use_urban_planning}





\cleardoublepage
 % \ref{chap:land_use_urban_planning}
    
% The next command tells RStudio to do "Compile PDF" on book.Rnw,
% instead of this chapter, thereby eliminating the need to switch back to book.Rnw 
% before making the book.
%!TEX root = ../../book.Rnw

\chapter{Government}
\label{chap:government}





\cleardoublepage
                     % \ref{chap:government}
    
% The next command tells RStudio to do "Compile PDF" on book.Rnw,
% instead of this chapter, thereby eliminating the need to switch back to book.Rnw 
% before making the book.
%!TEX root = ../../book.Rnw

\chapter{Personal action}
\label{chap:personal_action}





\cleardoublepage
             % \ref{chap:personal_action}

 %   \input{biblio}              % bibliography

    \bibliographystyle{unsrtnat}
    \bibliography{MCBook2021}   %chktex 11
    
    
    \cleardoublepage

 \backmatter                    % back matter
    
%blankpage

\chapter*{Author's Biography}
\markboth{AUTHOR'S BIOGRAPHY}{AUTHOR'S BIOGRAPHY}
\addcontentsline{toc}{chapter}{\protect\numberline{}{Author's Biography}}


\section*{Jeremy Van Antwerp}

% Adjust spacing so the photo looks nice in the paragraph.
\setlength{\intextsep}{-7pt}%
\setlength{\columnsep}{8pt}%
\begin{wrapfigure}{L}{0.25\textwidth}
  \begin{center}
    \includegraphics[width=0.25\textwidth]{headshots/jva-headshot.jpeg}
  \end{center}
\end{wrapfigure}
\textbf{Jeremy Van Antwerp} earned an M.S.\ and Ph.D.\ in chemical engineering 
from the University of Illinois at Urbana-Champaign. 
He is Professor of Engineering (chemical concentration)
at Calvin University in Grand Rapids, MI, USA
and has been a research affiliate and visiting scholar at the Massachusetts 
Institute of Technology, a journal editor, and a consultant in the pharmaceutical 
industry.
His research area is systems and process control.
He holds one US patent and is the author of \emph{Identifcation and control of 
sheet and film processes}~\cite{FeaVB2000}.
His poetry, cartoons, and puzzles have appeared in \emph{IEEE Control Systems}.
He is passionate about excellence and high-impact practices in education.
In 2001 he was awarded a Michigan Traditional Arts apprenticeship grant.


\section*{Matthew Kuperus Heun}

% Adjust spacing so the photo looks nice in the paragraph.
\setlength{\intextsep}{-7pt}%
\setlength{\columnsep}{8pt}%
\begin{wrapfigure}{L}{0.25\textwidth}
  \begin{center}
    \includegraphics[width=0.25\textwidth]{headshots/Heun-pub-photo-headshot}
  \end{center}
\end{wrapfigure}
\textbf{Matthew Kuperus Heun} is Professor of Engineering 
(mechanical concentration)
at Calvin University in Grand Rapids, MI, USA.
He earned an M.S.\ and Ph.D.\ in mechanical engineering from 
the University of Illinois at Urbana-Champaign and 
later worked at NASA's Jet Propulsion Laboratory and at Global Aerospace Corporation. 
He has been a visiting scholar at the Centre for Renewable and Sustainable Energy Studies 
at the University of Stellenbosch, South Africa. 
His long-term research question is 
``What is the relationship between energy and the economy when viewed through the lens of sustainability?''
In addition to scores of articles, he is lead author of 
\emph{Beyond GDP: National accounting in the age of resource depletion}~\citep{Heun:2015aa} 
and a co-editor of
\emph{Beyond Stewardship: New approaches to creation care}~\citep{Warners:2019aa}.
                 % author(s)'s short bio(s)
								% please see the bios folder for some pdf examples
    
    \cleardoublepage

    
% please use TexLive 2014 or later with the M&C macros freely
% available from tug.org or use any other recent version of LaTeX

\documentclass{book}\usepackage[]{graphicx}\usepackage[]{color}
% maxwidth is the original width if it is less than linewidth
% otherwise use linewidth (to make sure the graphics do not exceed the margin)
\makeatletter
\def\maxwidth{ %
  \ifdim\Gin@nat@width>\linewidth
    \linewidth
  \else
    \Gin@nat@width
  \fi
}
\makeatother

\definecolor{fgcolor}{rgb}{0.345, 0.345, 0.345}
\newcommand{\hlnum}[1]{\textcolor[rgb]{0.686,0.059,0.569}{#1}}%
\newcommand{\hlstr}[1]{\textcolor[rgb]{0.192,0.494,0.8}{#1}}%
\newcommand{\hlcom}[1]{\textcolor[rgb]{0.678,0.584,0.686}{\textit{#1}}}%
\newcommand{\hlopt}[1]{\textcolor[rgb]{0,0,0}{#1}}%
\newcommand{\hlstd}[1]{\textcolor[rgb]{0.345,0.345,0.345}{#1}}%
\newcommand{\hlkwa}[1]{\textcolor[rgb]{0.161,0.373,0.58}{\textbf{#1}}}%
\newcommand{\hlkwb}[1]{\textcolor[rgb]{0.69,0.353,0.396}{#1}}%
\newcommand{\hlkwc}[1]{\textcolor[rgb]{0.333,0.667,0.333}{#1}}%
\newcommand{\hlkwd}[1]{\textcolor[rgb]{0.737,0.353,0.396}{\textbf{#1}}}%
\let\hlipl\hlkwb

\usepackage{framed}
\makeatletter
\newenvironment{kframe}{%
 \def\at@end@of@kframe{}%
 \ifinner\ifhmode%
  \def\at@end@of@kframe{\end{minipage}}%
  \begin{minipage}{\columnwidth}%
 \fi\fi%
 \def\FrameCommand##1{\hskip\@totalleftmargin \hskip-\fboxsep
 \colorbox{shadecolor}{##1}\hskip-\fboxsep
     % There is no \\@totalrightmargin, so:
     \hskip-\linewidth \hskip-\@totalleftmargin \hskip\columnwidth}%
 \MakeFramed {\advance\hsize-\width
   \@totalleftmargin\z@ \linewidth\hsize
   \@setminipage}}%
 {\par\unskip\endMakeFramed%
 \at@end@of@kframe}
\makeatother

\definecolor{shadecolor}{rgb}{.97, .97, .97}
\definecolor{messagecolor}{rgb}{0, 0, 0}
\definecolor{warningcolor}{rgb}{1, 0, 1}
\definecolor{errorcolor}{rgb}{1, 0, 0}
\newenvironment{knitrout}{}{} % an empty environment to be redefined in TeX

\usepackage{alltt}

%the main style; default LibreCaslon font
\usepackage[raggedsec]{morgan2}
\usepackage{morgan-defs}

%to use Times New Roman, instead of LibreCaslon, please uncomment the next line
%\morgansetup{fontsetup=times}

% bibliography
% \usepackage[comma,sort,authoryear]{natbib}         % author-year
\usepackage[comma,sort,numbers]{natbib}            % numbered

% Attempting per-chapter bibliographies. Not working yet.
% \usepackage[sectionbib,compress]{natbib}% provides a bibliography for each chapter
% \usepackage{chapterbib}


% 
% Begin LaTeX packages imported by the authors.
% 

\usepackage{booktabs}    % For awesome table formatting
\usepackage[inline]{enumitem}  % For inline enumerate* lists
\usepackage{gensymb}     % For the \degree command
\usepackage{microtype}   % For (more) beautiful typesetting.
\usepackage{rotating}    % For rotating figures.
\usepackage{wrapfig}     % To wrap text around figures.

% 
% Begin LaTeX macros created by the authors.
% 

\newcommand{\degC}{\degree C}

% 
% Begin R and knitr setup
% 






%
% End author additions
% 


\setcounter{secnumdepth}{2}

\graphicspath{{./figures/}}     % folder for the figures in your book

\PassOptionsToPackage{hyphens}{url}
\usepackage[colorlinks=true,linkcolor=MyDarkBlue,
citecolor=MyDarkBlue,filecolor=MyDarkBlue,urlcolor=MyDarkBlue]{hyperref}

\renewcommand{\UrlBreaks}{\do\.\do\@\do\\\do\/\do\!\do\_\do\|\do\;\do\>\do\]%
\do\)\do\,\do\?\do\&\do\'\do+\do\=\do\#%
\do\a\do\b\do\c\do\d\do\e\do\f\do\g\do\h\do\i\do\j%
\do\k\do\l\do\m\do\n\do\o\do\p\do\q\do\r\do\s\do\t%
\do\u\do\v\do\w\do\x\do\y\do\z\do\A\do\B\do\C\do\D%
\do\E\do\F\do\G\do\H\do\I\do\J\do\K\do\L\do\M\do\N%
\do\O\do\P\do\Q\do\R\do\S\do\T\do\U\do\V\do\W\do\X%
\do\Y\do\Z}
\makeatletter
\g@addto@macro{\UrlBreaks}{\UrlOrds}
\renewcommand{\ALG@name}{\color{black}Algorithm}
\makeatother



\makeindex{}                     % if you are creating an index for your book
\IfFileExists{upquote.sty}{\usepackage{upquote}}{}
\begin{document}
% \SweaveOpts{concordance=TRUE}  % Commented on 17 Oct 2019 to eliminate some errors.

 \frontmatter                   % we'll produce the fm for you; 
    
\def\HALFTITLE{A Framework for Sustainability\\ Thinking: \\%The Multifaceted Nature\\
		%	of Sustainability Challenges:\\
			An Engineering Perspective \\ 
			\textcolor{red}{version $\alpha$ \\}
			\textcolor{black}{\today}}
\def\TITLE{A Framework for Sustainability Thinking \\ %The Multifaceted Nature of Sustainability Challenges\\ 
                        \textcolor{red}{version $\alpha$} \\ \textcolor{black}{\today}}
\def\AUTHORA{Jeremy Van Antwerp}
\def\AFFILIATIONA{Engineering Department, Calvin University}
\def\AUTHORB{Matthew Kuperus Heun}
\def\AFFILIATIONB{Engineering Department, Calvin University}
\def\AUTHORS{\AUTHORA\ and \AUTHORB}
\def\LECTURE{\ \#4}
\def\lcSYNTHESIS{Synthesis Lectures on Sustainable Development}
\def\SYNTHESIS{\MakeUppercase{\textit{\lcSYNTHESIS}}}

\def\EDITOR{xxxx, \textit{yyyy}}


%%%%%%%%%%%% HALF TITLE PAGE 1-2

\thispagestyle{emptyrule}

\halftitle{\HALFTITLE}

\clearpage


%%%%%%%%%%%% FULL TITLE 5

\blankpage

\thispagestyle{emptyrule}
\title{\TITLE}

\vspace*{2pc}
\authorname{\AUTHORA}
\authoraffiliation{\AFFILIATIONA}

\vspace*{1pc}
\authorname{\AUTHORB}
\authoraffiliation{\AFFILIATIONB}

\vfill
\synthesis{\SYNTHESIS\LECTURE}
\morganlogo

\clearpage


%%%%%%%%%%%% ABSTRACT AND KEYWORDS   6

\thispagestyle{emptyrule}

\ABSTRACT
\noindent
%The abstract goes here. 
%The Abstract and the keywords have to fit in this page.
This book is intended to be an introduction to the many challenges of 
sustainability.
It provides basic facts, figures, and information related to sustainability
in a way that is easy to digest.
The focus is on a quantitative understanding of the challenges of sustainability.
The first half of the book develops a framework for sustainability thinking.
The second half considers application areas and personal and corporate responses
to sustainability challenges.
Throughout, the end-of-chapter discussion questions focus on tradeoffs among 
competing goods and the ethical and social implications of decisions related to
sustainability.
This book was written for a university seminar course on sustainability but could
be used in other small-group discussion formats. 
It is intended to be easy to read but hard to digest.
%end ABSTRACT

\keywords{% %% take a look in the index to see what should get added as keywords!
%xxx, yyyy, zz
Sustainability, sustainable development, systems modeling
}

\vfill

\clearpage


%%%%%%%%%%%%% DEDICATION  7-8


%{
%\clearpage
%\thispagestyle{plain}
%
%\vspace*{13pc}\Large\it
%\centerline{To Eric, Jacob, and my parents.}
%
%}
%
%\clearpage
 

%%%%%%%%%%%%%%%%%%%%%%%% TOC %%%%%%%%%%%%%%%%%%

%blankpage

{
\pagestyle{plain}
\tableofcontents
}

\clearpage
 
 
%%%%%%%%%%%%%%%%%%%%%%%% PREFACE %%%%%%%%%%%%%%%%%%%
 
 
\blankpage

{
\chapter*{Preface}
\addcontentsline{toc}{chapter}{\protect\numberline{}{Preface}}
\thispagestyle{plain}
\markboth{PREFACE}{PREFACE}

\noindent
\section*{Who should use this book} 
This book is for anyone who is interested in the many challenges of 
sustainability.
It is an introduction and, as such, does not assume any prior background knowledge,
apart from just a little bit of math.
The book is not intended to be comprehensive but is intended to be easy to read.\\

\section*{How to use this book} 
Read it. Seriously; that's what you do with books.
This book was written to support a one semester-hour university seminar class
on sustainability. The primary focus of class time is in-class discussion, so 
the end-of-chapter discussion questions should provide significant value. 
Therefore, individual chapters, or sets of chapters, can easily be
used to support a wide variety of class structures, types, and levels.
Students should read the text on their own, out of class, to get 
basic information related to  the challenges of sustainability. 
We encourage instructors to use a basic and straightforward quiz on the reading 
through your course management platform to hold students accountable for the 
reading. 
Students should review discussion questions ahead of each class and 
prepare for discussion. \\

\section*{What you should get out of this book} 
This book contains four main contributions.
\begin{itemize}
\item Basic facts, figures, and information related to sustainability. 
We don't assume you start with any prior knowledge but by the time you finish 
reading this book you will know basic facts like carbon dioxide 
concentration in the atmosphere and what are the anthropogenic sources of 
carbon dioxide. 
Furthermore, you should have a \emph{sense of scale} related to sustainability 
information. 
That is, which things are big and important and which are smaller and less relavant. 
To this end, the information is presented \emph{graphically} as much as possible 
to show relative magnitudes.
\item We provide an explanatory framework or scaffold for thinking about sustainability,
primarily in chapters \ref{chap:introduction}-\ref{chap:carbon_intensity}.
\item The end-of-chapter discussion questions get at moral, ethical, philosophical,
and practical aspects of sustainability. 
They don't have %\sout{right} 
objective answers. 
Instead, the questions should help illustrate the role that worldview, values, 
preferences, and \emph{a priori} assumptions have in how we approach sustainability. 
%social pillar?
\item The end-of-chapter project questions range in difficulty from a long
homework problem to a graduate thesis project. 
If graduate study on a sustainability topic is in your future, perhaps you can 
take inspiration from one or more of these project questions. 
Instructors can use one -- or perhaps more -- of these questions for a 
semester-long project, optionally with in-class presentations, for a class.
\end{itemize}
By the end of the book you should see that sustainability is an important, 
urgent, difficult, but solvable problem that requires personal and corporate 
commitments. \\

\section*{Text, organization, and schedule} 
At one chapter a week, the twelve chapters in this book don't quite fill a typical 
university semester. 
Of course, you don't have to be in a university context to use this book. 
A book club or church small group are also appropriate audiences -- sustainability 
is important for everyone. 
However, if you are in a university context, here are a few ways that you could fill the 
remaining weeks of your semester. 
Use class time for project presentations, as mentioned above. 
Perhaps groups of students could each be assigned different project questions and 
at the end of the semester, each could share their results.
Alternatively, a class session or sessions could be devoted to a topic or topics 
your instructor finds personally interesting, like biofuels, grid-scale storage, 
or media coverage of sustainability. 
At Calvin University, our class devotes a week to worldview implications for 
sustainability. 
You could use all or part(s) of \emph{Beyond stewardship: New approaches to 
creation care}~\citep{Warners:2019aa}
or our paper \emph{Current disciplines and worldviews are insufficient
to address sustainability challenges} (2019) as a starting place for discussion. 
% CES website doesn't have 2019 proceedings!
On the other hand, instructors don't need to use all of the book. 
IEEE Xplore makes it easy to assign either a single chapter or a group of chapters 
as supplementary material for many different types of classes.
Lastly, if you read all the way to the end of the Preface, great job! 
Give youself a gold star. 
Now, on to the good stuff...


\vspace*{2pc}
\noindent\AUTHORS\\
\noindent December 2021
}

\clearpage


%%%%%%%%%%%%%%%%%%%%%%%% ACK %%%%%%%%%%%%%%%%%%%


\blankpage

\chapter*{Acknowledgments}
\addcontentsline{toc}{chapter}{\protect\numberline{}{Acknowledgments}}
\thispagestyle{plain}
\markboth{ACKNOWLEDGMENTS}{ACKNOWLEDGMENTS}

\noindent
Thanks to the students of the Sustainability Challenges course and our colleague 
Julie Wildschut for teaching the class in fall 2021.

Thanks to research assistants Larisa Tomeci, Henos Tadesse, and Joshua Broekhuisen
for help with sourcing data and creating figures for the text. 
% was there a 4th student?

Thank you to the many reviewers who provided feedback on early versions to the book:
James VanAntwerp, Julie Wildschut, Becky Haney, James Skillen, Uko Zylstra.

% Jeff Jewett, Shannon Savage, Leonides Murembya, Steve McMullen, Nathan Grawe
% Jennifer VanAntwerp,  Glen VanAntwerp, ...
% Derek Schuurman, Sa BA, Tom Ackerman

Thanks to the Calvin University Board of Trustees and Joel and Linda Zylstra
for funding to create and revise the text.
Thanks also to Tim Koning of Liquid Haulers Maintenance for funding on biofuels research.



\vspace*{2pc}
\noindent\AUTHORS\\
\noindent December 2021
 
\clearpage

\blankpage

                  % we'll need an Abstract and Keywords
								% please see the abs-pref folder for some pdf examples

 \mainmatter
    
% The next command tells RStudio to do "Compile PDF" on book.Rnw,
% instead of this chapter, thereby eliminating the need to switch back to book.Rnw 
% before making the book.
%!TEX root = ../../book.Rnw

\chapter{Introduction}
\label{chap:introduction}

Let’s begin with a story. Once upon a time, a king was challenged to a game of chess by a poor farmer. If he lost, the king would have to pay the farmer a chessboard’s worth of grain, defined as one kernel of grain on the first square, two grains on the second square, four grains on the fourth square, and so on, with each subsequent square receiving double what was on the previous square. The king accepted the wager but, because this is fable, lost the game. A chessboard has eight rows and eight columns, so the king had to pay 264 (less one) grains of rice. Since it takes seven thousand grains of rice to make a pound (~65 mg each), the king owed her 1.195 exograms (one exogram is 1015 kg), which is about 12\% of the weight of everything ever produced by the human species [1]. At US\$6.85/bushel, the grain would be worth about 300 trillion dollars, which is (roughly) ten times more money than exists in the world. 
%Clean up. Maybe kg/person in the world and $/person. 

This story, which goes back at least until the 1200s, illustrates several aspects of the sustainability problem that have been evident for centuries.  Exponential growth is not sustainable. 
Big numbers are hard to grok. % Cognitive biases Decartes’ chiliagon Numerical cognition Subitizing Dunbar number 50-year time frame.
It’s not OK to take giant risks.

The purpose of this book is to summarize the evidence that humans are not living sustainably and to show the characteristics of sustainable societies. The choices and paths to sustainability are bristling with value judgements and moral choices. This book attempts to provide a coherent framework for discussing sustainability that is grounded in a sense of scale. Each chapter is short, with focused discussion questions at the end of each chapter. The questions, which mostly focus on tradeoffs, do not have “right answers,” instead, different answers are indicative of different preferences. Basic data and key concepts start each chapter. This book is not comprehensive; the intent is to equip readers with basic knowledge (informed by scale) so we can grapple with tough moral questions. We intend the book to be easy to read but hard to digest. 

The remainder of this chapter will summarize key themes of the book. Sustainability is existentially important but very difficult for humans to achieve, for several reasons. The numbers and timescales involved are difficult for humans to grasp for a variety of reasons. Sustainability challenges are complex because they involved a variety of interrelated problems from different domain areas. The challenges associated with transitioning to a sustainable existence are complex and interconnected. It is the intention of this book to explore the challenges of sustainability. While we may, at times, point to directions for improved sustainability, it is beyond our scope (and indeed our ability) to provide solutions for all sustainability problems. 

\section{What is sustainability?}
Are humans living "sustainably?" If not, we will -- at some point -- either cease to exist as a species or (at least) experience drastic reductions in our standard of living. Thus, the question of sustainability is of utmost importance. There are many answers to this question. To the novice, many answers may seem to indicate lack of agreement. However, sustainability is almost a self-defining concept. Different answers indicate different assumptions, different priorities, and different boundaries (that is, what is the system under consideration). 

When considering the meaning of sustainability, the two most important questions are "sustaining what?" and "for how long?" The second of these questions is perhaps easier to answer. While humans seem to have difficulty planning for time scales significantly longer than the human lifespan, true sustainability is achieved only if the answer is "indefinitely" or "forever."
The narrowest answer for what needs sustaining is human life and society, which perforce entails those ecosystem services necessary for human health and wellbeing. Beyond these basics, some people view the nonhuman world as having inherent worth or standing and, as such, to also be worth preserving, even if it has no (or negative, for example, smallpox) impact on human continuity. Unfortunately, humans don’t know clearly what pieces of the ecosystem are, in the long run, necessary for survival and which aren’t. For instance, would we be able to survive in a world without dandelions? Maybe. On the other hand, dandelions might be necessary for other organisms we depend on. Environmental science views the ecosystem in its entirety as a web, with all parts depending on all other parts. The ecosystem is viewed then, not as individual components with binary "needed" "not needed" classifications but as a whole that exists on a continuum from "fully functional" to "nonfunctional." The choices and paths to sustainability are bristling with value judgements and moral choices. How much risk do you think is appropriate? What do you value preserving? Different answers to these questions come from different a priori assumptions and values. 

The broadest definitions of sustainability include human-created devices and systems. Humanity could survive without the great pyramids of Giza, although we would have lost great cultural and historical artifacts. On the other hand, we may not survive if we don’t give up coal-fired power plants. Thus, the Three Pillars of sustainability are its environmental, economic, and social aspects. 
Environmental sustainability considers biophysical and thermodynamic constraints and includes issues such as pollution, resource depletion, habitat loss, and biodiversity. Economic sustainability involves questions of profit and loss, wealth management, and macroeconomic policy. Social sustainability comprises human and civil rights, suffering, and personal freedom. This book will discuss issues associated with each of these three areas of sustainability but with a focus on energy and carbon emissions because it is not the only sustainability problem, just the one that needs to be solved most urgently.

The most oft-quoted definition of sustainability comes from the Brundtland report "Sustainable development is development that meets the needs of the present without compromising the ability of future generations to meet their own needs." Note that sustainable development (an oxymoron?!) is not the same as sustainability and “needs” are subjective.

Herman Daly proposed the following three rules for achieving ecological (or environmental) sustainability:
\begin{enumerate}
\item Renewable resources such as fish, soil, and groundwater must be used no faster than the rate at which they regenerate.
\item Nonrenewable resources such as minerals and fossil fuels must be used no faster than renewable substitutes for them can be put into place.
\item Pollution and wastes must be emitted no faster than natural systems can absorb them, recycle them, or render them harmless.
\end{enumerate}

The core idea is that [the rate of] human impact on the ecosystem can be no greater than the ecosystem can absorb, which introduces the idea of the capacity of the ecosystem to provide resources and “take care of” wastes. We will examine the idea of capacity in greater detail in the next chapter.

\section{Impacts and their sources}
Two mathematical expressions, IPAT and Kaya have been used express human impact on the environment and its sources. IPAT expresses “impact” on the environment as the product of human population, affluence, and technology. The Kaya identity is a form of the IPAT equation but limits impact to CO2 emissions. Kaya also helpfully expresses the generic “technology” as the product of energy intensity of the economy and the carbon intensity of energy. 

The first section of this book is organized around a combination of the IPAT and Kaya approach. We keep the general “impact” of IPAT but use the more helpful expression of technology from Kaya. 
\begin{equation}\label{eq:IPARC}
Impact = Population \times Affluence \times Resource intensity \times Contamination intensity
\end{equation}
Impact is a vector quantity (that is, a list) including such things as global warming potential, aquifer depletion, and eutrophication potential. Population is the (scalar) number of people in the world. Affluence is the world per capita GDP. Resource intensity is the vector or list of resources necessary to produce one unit of world GDP. Lastly, impact intensity or contamination intensity is a matrix that lists the impacts of each type of resource.

IPARC is true because it is an identity, and it will provide a useful organizing framework for the following chapters. The terms are equilibrium values, that is they represent static, or steady-state levels, but do not give completely useful information about how changes in any one variable will affect the others. For instance, using resources more efficiently (improving resource intensity), does not lower impact. Instead, it leads to more affluence (see chapter 5). Likewise, it may not be possible to drive the impact(s) per unit of resource to zero because of diminishing returns on efficiency and tradeoffs that exist between different types of impacts.

\section{The mathematics of sustainability}
Returning to the wheat and chessboard problem, for a series of doubling numbers, like the number of grains on a chessboard square: $1, 2, 4, 8, 16, 32, …$, the sum of all the numbers in the series is always one less than the next number in the series. $1+2+4+8 = 15$, which is one less than $16$. $1+2+4+8+16 = 31$, which is one less than $32$. 

Extending the doubling on discrete squares or in discrete time to continuous time is exponential growth. Things that grow exponentially double in a fixed amount of time. An example with a negative growth rate is radioactive decay, where the corresponding time is called the half-life because the amount of time for half of the radioactivity to decay away is constant. When growth rate is constant, the amount of time to double is approximately 70 divided by the growth rate. With a fixed 3\% growth rate, GDP doubles every 23 years. If your investment portfolio has a 7\% return, your assets will double every 10 years. 

In the discrete case, like the wheat and chessboard, the next square is always one more than the sum of all the previous squares. In the continuous case, the amount in the next period (B) is always 1.71 times the sum of all previous amounts (A).
The result for sustainability is that the resources needed by the economy in the next 23 years – for instance, energy – will be more than the total resources (energy) used in all of history. Populations are modeled as having constant growth rates. We actively manage the economy have a constant growth rate (say a target of 3\% GDP growth per year). Because of (exponential) growth in population and affluence, humanity is not on a sustainable path. 

\section{CO2 and Energy}
The approach above (IPARC) is general in that all kinds of impacts can be considered. In this book we will consider many types of impacts and many kinds of sustainability (environmental, economic, and social). However, our focus – in the chapters – is mostly on energy use and CO2 emissions. Reasons for an energy/CO2 focus include the availability of data and that we feel it is (currently) the most urgent sustainability problem. End-of-chapter discussion questions tend to focus more on social aspects of sustainability and worldview issues.


\cleardoublepage
                 % \ref{chap:introduction}
    
% The next command tells RStudio to do "Compile PDF" on book.Rnw,
% instead of this chapter, thereby eliminating the need to switch back to book.Rnw 
% before making the book.
%!TEX root = ../../book.Rnw

\chapter{Planetary Boundaries}
\label{chap:planetary_boundaries}





\cleardoublepage
                     % \ref{chap:planetary_boundaries}
    
% The next command tells RStudio to do "Compile PDF" on book.Rnw,
% instead of this chapter, thereby eliminating the need to switch back to book.Rnw 
% before making the book.
%!TEX root = ../../book.Rnw

\chapter{Population}
\label{chap:population}





\cleardoublepage
                     % \ref{chap:population}
    
% The next command tells RStudio to do "Compile PDF" on book.Rnw,
% instead of this chapter, thereby eliminating the need to switch back to book.Rnw 
% before making the book.
%!TEX root = ../../book.Rnw

\chapter{Affluence}
\label{chap:affluence}





\cleardoublepage
                       % \ref{chap:affluence}
    
% The next command tells RStudio to do "Compile PDF" on book.Rnw,
% instead of this chapter, thereby eliminating the need to switch back to book.Rnw 
% before making the book.
%!TEX root = ../../book.Rnw

\chapter{Energy intensity of the economy}
\label{chap:energy_intensity}





\cleardoublepage
           % \ref{chap:energy_intensity}
    
% The next command tells RStudio to do "Compile PDF" on book.Rnw,
% instead of this chapter, thereby eliminating the need to switch back to book.Rnw 
% before making the book.
%!TEX root = ../../book.Rnw

\chapter{Carbon intensity of energy}
\label{chap:carbon_intensity}





\cleardoublepage
           % \ref{chap:chap:carbon_intensity}
    
% The next command tells RStudio to do "Compile PDF" on book.Rnw,
% instead of this chapter, thereby eliminating the need to switch back to book.Rnw 
% before making the book.
%!TEX root = ../../book.Rnw

\chapter{Households}
\label{chap:households}





\cleardoublepage
                     % \ref{chap:households}
    
% The next command tells RStudio to do "Compile PDF" on book.Rnw,
% instead of this chapter, thereby eliminating the need to switch back to book.Rnw 
% before making the book.
%!TEX root = ../../book.Rnw

\chapter{Transportation}
\label{chap:transportation}





\cleardoublepage
             % \ref{chap:transportation}
    
% The next command tells RStudio to do "Compile PDF" on book.Rnw,
% instead of this chapter, thereby eliminating the need to switch back to book.Rnw 
% before making the book.
%!TEX root = ../../book.Rnw

\chapter{Agriculture}
\label{chap:agriculture}





\cleardoublepage
                   % \ref{chap:agriculture}
    
% The next command tells RStudio to do "Compile PDF" on book.Rnw,
% instead of this chapter, thereby eliminating the need to switch back to book.Rnw 
% before making the book.
%!TEX root = ../../book.Rnw

\chapter{Land use and urban planning}
\label{chap:land_use_urban_planning}





\cleardoublepage
 % \ref{chap:land_use_urban_planning}
    
% The next command tells RStudio to do "Compile PDF" on book.Rnw,
% instead of this chapter, thereby eliminating the need to switch back to book.Rnw 
% before making the book.
%!TEX root = ../../book.Rnw

\chapter{Government}
\label{chap:government}





\cleardoublepage
                     % \ref{chap:government}
    
% The next command tells RStudio to do "Compile PDF" on book.Rnw,
% instead of this chapter, thereby eliminating the need to switch back to book.Rnw 
% before making the book.
%!TEX root = ../../book.Rnw

\chapter{Personal action}
\label{chap:personal_action}





\cleardoublepage
             % \ref{chap:personal_action}

 %   \input{biblio}              % bibliography

    \bibliographystyle{unsrtnat}
    \bibliography{MCBook2021}   %chktex 11
    
    
    \cleardoublepage

 \backmatter                    % back matter
    
%blankpage

\chapter*{Author's Biography}
\markboth{AUTHOR'S BIOGRAPHY}{AUTHOR'S BIOGRAPHY}
\addcontentsline{toc}{chapter}{\protect\numberline{}{Author's Biography}}


\section*{Jeremy Van Antwerp}

% Adjust spacing so the photo looks nice in the paragraph.
\setlength{\intextsep}{-7pt}%
\setlength{\columnsep}{8pt}%
\begin{wrapfigure}{L}{0.25\textwidth}
  \begin{center}
    \includegraphics[width=0.25\textwidth]{headshots/jva-headshot.jpeg}
  \end{center}
\end{wrapfigure}
\textbf{Jeremy Van Antwerp} earned an M.S.\ and Ph.D.\ in chemical engineering 
from the University of Illinois at Urbana-Champaign. 
He is Professor of Engineering (chemical concentration)
at Calvin University in Grand Rapids, MI, USA
and has been a research affiliate and visiting scholar at the Massachusetts 
Institute of Technology, a journal editor, and a consultant in the pharmaceutical 
industry.
His research area is systems and process control.
He holds one US patent and is the author of \emph{Identifcation and control of 
sheet and film processes}~\cite{FeaVB2000}.
His poetry, cartoons, and puzzles have appeared in \emph{IEEE Control Systems}.
He is passionate about excellence and high-impact practices in education.
In 2001 he was awarded a Michigan Traditional Arts apprenticeship grant.


\section*{Matthew Kuperus Heun}

% Adjust spacing so the photo looks nice in the paragraph.
\setlength{\intextsep}{-7pt}%
\setlength{\columnsep}{8pt}%
\begin{wrapfigure}{L}{0.25\textwidth}
  \begin{center}
    \includegraphics[width=0.25\textwidth]{headshots/Heun-pub-photo-headshot}
  \end{center}
\end{wrapfigure}
\textbf{Matthew Kuperus Heun} is Professor of Engineering 
(mechanical concentration)
at Calvin University in Grand Rapids, MI, USA.
He earned an M.S.\ and Ph.D.\ in mechanical engineering from 
the University of Illinois at Urbana-Champaign and 
later worked at NASA's Jet Propulsion Laboratory and at Global Aerospace Corporation. 
He has been a visiting scholar at the Centre for Renewable and Sustainable Energy Studies 
at the University of Stellenbosch, South Africa. 
His long-term research question is 
``What is the relationship between energy and the economy when viewed through the lens of sustainability?''
In addition to scores of articles, he is lead author of 
\emph{Beyond GDP: National accounting in the age of resource depletion}~\citep{Heun:2015aa} 
and a co-editor of
\emph{Beyond Stewardship: New approaches to creation care}~\citep{Warners:2019aa}.
                 % author(s)'s short bio(s)
								% please see the bios folder for some pdf examples
    
    \cleardoublepage

    
% please use TexLive 2014 or later with the M&C macros freely
% available from tug.org or use any other recent version of LaTeX

\documentclass{book}\usepackage[]{graphicx}\usepackage[]{color}
% maxwidth is the original width if it is less than linewidth
% otherwise use linewidth (to make sure the graphics do not exceed the margin)
\makeatletter
\def\maxwidth{ %
  \ifdim\Gin@nat@width>\linewidth
    \linewidth
  \else
    \Gin@nat@width
  \fi
}
\makeatother

\definecolor{fgcolor}{rgb}{0.345, 0.345, 0.345}
\newcommand{\hlnum}[1]{\textcolor[rgb]{0.686,0.059,0.569}{#1}}%
\newcommand{\hlstr}[1]{\textcolor[rgb]{0.192,0.494,0.8}{#1}}%
\newcommand{\hlcom}[1]{\textcolor[rgb]{0.678,0.584,0.686}{\textit{#1}}}%
\newcommand{\hlopt}[1]{\textcolor[rgb]{0,0,0}{#1}}%
\newcommand{\hlstd}[1]{\textcolor[rgb]{0.345,0.345,0.345}{#1}}%
\newcommand{\hlkwa}[1]{\textcolor[rgb]{0.161,0.373,0.58}{\textbf{#1}}}%
\newcommand{\hlkwb}[1]{\textcolor[rgb]{0.69,0.353,0.396}{#1}}%
\newcommand{\hlkwc}[1]{\textcolor[rgb]{0.333,0.667,0.333}{#1}}%
\newcommand{\hlkwd}[1]{\textcolor[rgb]{0.737,0.353,0.396}{\textbf{#1}}}%
\let\hlipl\hlkwb

\usepackage{framed}
\makeatletter
\newenvironment{kframe}{%
 \def\at@end@of@kframe{}%
 \ifinner\ifhmode%
  \def\at@end@of@kframe{\end{minipage}}%
  \begin{minipage}{\columnwidth}%
 \fi\fi%
 \def\FrameCommand##1{\hskip\@totalleftmargin \hskip-\fboxsep
 \colorbox{shadecolor}{##1}\hskip-\fboxsep
     % There is no \\@totalrightmargin, so:
     \hskip-\linewidth \hskip-\@totalleftmargin \hskip\columnwidth}%
 \MakeFramed {\advance\hsize-\width
   \@totalleftmargin\z@ \linewidth\hsize
   \@setminipage}}%
 {\par\unskip\endMakeFramed%
 \at@end@of@kframe}
\makeatother

\definecolor{shadecolor}{rgb}{.97, .97, .97}
\definecolor{messagecolor}{rgb}{0, 0, 0}
\definecolor{warningcolor}{rgb}{1, 0, 1}
\definecolor{errorcolor}{rgb}{1, 0, 0}
\newenvironment{knitrout}{}{} % an empty environment to be redefined in TeX

\usepackage{alltt}

%the main style; default LibreCaslon font
\usepackage[raggedsec]{morgan2}
\usepackage{morgan-defs}

%to use Times New Roman, instead of LibreCaslon, please uncomment the next line
%\morgansetup{fontsetup=times}

% bibliography
% \usepackage[comma,sort,authoryear]{natbib}         % author-year
\usepackage[comma,sort,numbers]{natbib}            % numbered

% Attempting per-chapter bibliographies. Not working yet.
% \usepackage[sectionbib,compress]{natbib}% provides a bibliography for each chapter
% \usepackage{chapterbib}


% 
% Begin LaTeX packages imported by the authors.
% 

\usepackage{booktabs}    % For awesome table formatting
\usepackage[inline]{enumitem}  % For inline enumerate* lists
\usepackage{gensymb}     % For the \degree command
\usepackage{microtype}   % For (more) beautiful typesetting.
\usepackage{rotating}    % For rotating figures.
\usepackage{wrapfig}     % To wrap text around figures.

% 
% Begin LaTeX macros created by the authors.
% 

\newcommand{\degC}{\degree C}

% 
% Begin R and knitr setup
% 






%
% End author additions
% 


\setcounter{secnumdepth}{2}

\graphicspath{{./figures/}}     % folder for the figures in your book

\PassOptionsToPackage{hyphens}{url}
\usepackage[colorlinks=true,linkcolor=MyDarkBlue,
citecolor=MyDarkBlue,filecolor=MyDarkBlue,urlcolor=MyDarkBlue]{hyperref}

\renewcommand{\UrlBreaks}{\do\.\do\@\do\\\do\/\do\!\do\_\do\|\do\;\do\>\do\]%
\do\)\do\,\do\?\do\&\do\'\do+\do\=\do\#%
\do\a\do\b\do\c\do\d\do\e\do\f\do\g\do\h\do\i\do\j%
\do\k\do\l\do\m\do\n\do\o\do\p\do\q\do\r\do\s\do\t%
\do\u\do\v\do\w\do\x\do\y\do\z\do\A\do\B\do\C\do\D%
\do\E\do\F\do\G\do\H\do\I\do\J\do\K\do\L\do\M\do\N%
\do\O\do\P\do\Q\do\R\do\S\do\T\do\U\do\V\do\W\do\X%
\do\Y\do\Z}
\makeatletter
\g@addto@macro{\UrlBreaks}{\UrlOrds}
\renewcommand{\ALG@name}{\color{black}Algorithm}
\makeatother



\makeindex{}                     % if you are creating an index for your book
\IfFileExists{upquote.sty}{\usepackage{upquote}}{}
\begin{document}
% \SweaveOpts{concordance=TRUE}  % Commented on 17 Oct 2019 to eliminate some errors.

 \frontmatter                   % we'll produce the fm for you; 
    \input{fm}                  % we'll need an Abstract and Keywords
								% please see the abs-pref folder for some pdf examples

 \mainmatter
    
% The next command tells RStudio to do "Compile PDF" on book.Rnw,
% instead of this chapter, thereby eliminating the need to switch back to book.Rnw 
% before making the book.
%!TEX root = ../../book.Rnw

\chapter{Introduction}
\label{chap:introduction}

Let’s begin with a story. Once upon a time, a king was challenged to a game of chess by a poor farmer. If he lost, the king would have to pay the farmer a chessboard’s worth of grain, defined as one kernel of grain on the first square, two grains on the second square, four grains on the fourth square, and so on, with each subsequent square receiving double what was on the previous square. The king accepted the wager but, because this is fable, lost the game. A chessboard has eight rows and eight columns, so the king had to pay 264 (less one) grains of rice. Since it takes seven thousand grains of rice to make a pound (~65 mg each), the king owed her 1.195 exograms (one exogram is 1015 kg), which is about 12\% of the weight of everything ever produced by the human species [1]. At US\$6.85/bushel, the grain would be worth about 300 trillion dollars, which is (roughly) ten times more money than exists in the world. 
%Clean up. Maybe kg/person in the world and $/person. 

This story, which goes back at least until the 1200s, illustrates several aspects of the sustainability problem that have been evident for centuries.  Exponential growth is not sustainable. 
Big numbers are hard to grok. % Cognitive biases Decartes’ chiliagon Numerical cognition Subitizing Dunbar number 50-year time frame.
It’s not OK to take giant risks.

The purpose of this book is to summarize the evidence that humans are not living sustainably and to show the characteristics of sustainable societies. The choices and paths to sustainability are bristling with value judgements and moral choices. This book attempts to provide a coherent framework for discussing sustainability that is grounded in a sense of scale. Each chapter is short, with focused discussion questions at the end of each chapter. The questions, which mostly focus on tradeoffs, do not have “right answers,” instead, different answers are indicative of different preferences. Basic data and key concepts start each chapter. This book is not comprehensive; the intent is to equip readers with basic knowledge (informed by scale) so we can grapple with tough moral questions. We intend the book to be easy to read but hard to digest. 

The remainder of this chapter will summarize key themes of the book. Sustainability is existentially important but very difficult for humans to achieve, for several reasons. The numbers and timescales involved are difficult for humans to grasp for a variety of reasons. Sustainability challenges are complex because they involved a variety of interrelated problems from different domain areas. The challenges associated with transitioning to a sustainable existence are complex and interconnected. It is the intention of this book to explore the challenges of sustainability. While we may, at times, point to directions for improved sustainability, it is beyond our scope (and indeed our ability) to provide solutions for all sustainability problems. 

\section{What is sustainability?}
Are humans living "sustainably?" If not, we will -- at some point -- either cease to exist as a species or (at least) experience drastic reductions in our standard of living. Thus, the question of sustainability is of utmost importance. There are many answers to this question. To the novice, many answers may seem to indicate lack of agreement. However, sustainability is almost a self-defining concept. Different answers indicate different assumptions, different priorities, and different boundaries (that is, what is the system under consideration). 

When considering the meaning of sustainability, the two most important questions are "sustaining what?" and "for how long?" The second of these questions is perhaps easier to answer. While humans seem to have difficulty planning for time scales significantly longer than the human lifespan, true sustainability is achieved only if the answer is "indefinitely" or "forever."
The narrowest answer for what needs sustaining is human life and society, which perforce entails those ecosystem services necessary for human health and wellbeing. Beyond these basics, some people view the nonhuman world as having inherent worth or standing and, as such, to also be worth preserving, even if it has no (or negative, for example, smallpox) impact on human continuity. Unfortunately, humans don’t know clearly what pieces of the ecosystem are, in the long run, necessary for survival and which aren’t. For instance, would we be able to survive in a world without dandelions? Maybe. On the other hand, dandelions might be necessary for other organisms we depend on. Environmental science views the ecosystem in its entirety as a web, with all parts depending on all other parts. The ecosystem is viewed then, not as individual components with binary "needed" "not needed" classifications but as a whole that exists on a continuum from "fully functional" to "nonfunctional." The choices and paths to sustainability are bristling with value judgements and moral choices. How much risk do you think is appropriate? What do you value preserving? Different answers to these questions come from different a priori assumptions and values. 

The broadest definitions of sustainability include human-created devices and systems. Humanity could survive without the great pyramids of Giza, although we would have lost great cultural and historical artifacts. On the other hand, we may not survive if we don’t give up coal-fired power plants. Thus, the Three Pillars of sustainability are its environmental, economic, and social aspects. 
Environmental sustainability considers biophysical and thermodynamic constraints and includes issues such as pollution, resource depletion, habitat loss, and biodiversity. Economic sustainability involves questions of profit and loss, wealth management, and macroeconomic policy. Social sustainability comprises human and civil rights, suffering, and personal freedom. This book will discuss issues associated with each of these three areas of sustainability but with a focus on energy and carbon emissions because it is not the only sustainability problem, just the one that needs to be solved most urgently.

The most oft-quoted definition of sustainability comes from the Brundtland report "Sustainable development is development that meets the needs of the present without compromising the ability of future generations to meet their own needs." Note that sustainable development (an oxymoron?!) is not the same as sustainability and “needs” are subjective.

Herman Daly proposed the following three rules for achieving ecological (or environmental) sustainability:
\begin{enumerate}
\item Renewable resources such as fish, soil, and groundwater must be used no faster than the rate at which they regenerate.
\item Nonrenewable resources such as minerals and fossil fuels must be used no faster than renewable substitutes for them can be put into place.
\item Pollution and wastes must be emitted no faster than natural systems can absorb them, recycle them, or render them harmless.
\end{enumerate}

The core idea is that [the rate of] human impact on the ecosystem can be no greater than the ecosystem can absorb, which introduces the idea of the capacity of the ecosystem to provide resources and “take care of” wastes. We will examine the idea of capacity in greater detail in the next chapter.

\section{Impacts and their sources}
Two mathematical expressions, IPAT and Kaya have been used express human impact on the environment and its sources. IPAT expresses “impact” on the environment as the product of human population, affluence, and technology. The Kaya identity is a form of the IPAT equation but limits impact to CO2 emissions. Kaya also helpfully expresses the generic “technology” as the product of energy intensity of the economy and the carbon intensity of energy. 

The first section of this book is organized around a combination of the IPAT and Kaya approach. We keep the general “impact” of IPAT but use the more helpful expression of technology from Kaya. 
\begin{equation}\label{eq:IPARC}
Impact = Population \times Affluence \times Resource intensity \times Contamination intensity
\end{equation}
Impact is a vector quantity (that is, a list) including such things as global warming potential, aquifer depletion, and eutrophication potential. Population is the (scalar) number of people in the world. Affluence is the world per capita GDP. Resource intensity is the vector or list of resources necessary to produce one unit of world GDP. Lastly, impact intensity or contamination intensity is a matrix that lists the impacts of each type of resource.

IPARC is true because it is an identity, and it will provide a useful organizing framework for the following chapters. The terms are equilibrium values, that is they represent static, or steady-state levels, but do not give completely useful information about how changes in any one variable will affect the others. For instance, using resources more efficiently (improving resource intensity), does not lower impact. Instead, it leads to more affluence (see chapter 5). Likewise, it may not be possible to drive the impact(s) per unit of resource to zero because of diminishing returns on efficiency and tradeoffs that exist between different types of impacts.

\section{The mathematics of sustainability}
Returning to the wheat and chessboard problem, for a series of doubling numbers, like the number of grains on a chessboard square: $1, 2, 4, 8, 16, 32, …$, the sum of all the numbers in the series is always one less than the next number in the series. $1+2+4+8 = 15$, which is one less than $16$. $1+2+4+8+16 = 31$, which is one less than $32$. 

Extending the doubling on discrete squares or in discrete time to continuous time is exponential growth. Things that grow exponentially double in a fixed amount of time. An example with a negative growth rate is radioactive decay, where the corresponding time is called the half-life because the amount of time for half of the radioactivity to decay away is constant. When growth rate is constant, the amount of time to double is approximately 70 divided by the growth rate. With a fixed 3\% growth rate, GDP doubles every 23 years. If your investment portfolio has a 7\% return, your assets will double every 10 years. 

In the discrete case, like the wheat and chessboard, the next square is always one more than the sum of all the previous squares. In the continuous case, the amount in the next period (B) is always 1.71 times the sum of all previous amounts (A).
The result for sustainability is that the resources needed by the economy in the next 23 years – for instance, energy – will be more than the total resources (energy) used in all of history. Populations are modeled as having constant growth rates. We actively manage the economy have a constant growth rate (say a target of 3\% GDP growth per year). Because of (exponential) growth in population and affluence, humanity is not on a sustainable path. 

\section{CO2 and Energy}
The approach above (IPARC) is general in that all kinds of impacts can be considered. In this book we will consider many types of impacts and many kinds of sustainability (environmental, economic, and social). However, our focus – in the chapters – is mostly on energy use and CO2 emissions. Reasons for an energy/CO2 focus include the availability of data and that we feel it is (currently) the most urgent sustainability problem. End-of-chapter discussion questions tend to focus more on social aspects of sustainability and worldview issues.


\cleardoublepage
                 % \ref{chap:introduction}
    
% The next command tells RStudio to do "Compile PDF" on book.Rnw,
% instead of this chapter, thereby eliminating the need to switch back to book.Rnw 
% before making the book.
%!TEX root = ../../book.Rnw

\chapter{Planetary Boundaries}
\label{chap:planetary_boundaries}





\cleardoublepage
                     % \ref{chap:planetary_boundaries}
    
% The next command tells RStudio to do "Compile PDF" on book.Rnw,
% instead of this chapter, thereby eliminating the need to switch back to book.Rnw 
% before making the book.
%!TEX root = ../../book.Rnw

\chapter{Population}
\label{chap:population}





\cleardoublepage
                     % \ref{chap:population}
    
% The next command tells RStudio to do "Compile PDF" on book.Rnw,
% instead of this chapter, thereby eliminating the need to switch back to book.Rnw 
% before making the book.
%!TEX root = ../../book.Rnw

\chapter{Affluence}
\label{chap:affluence}





\cleardoublepage
                       % \ref{chap:affluence}
    
% The next command tells RStudio to do "Compile PDF" on book.Rnw,
% instead of this chapter, thereby eliminating the need to switch back to book.Rnw 
% before making the book.
%!TEX root = ../../book.Rnw

\chapter{Energy intensity of the economy}
\label{chap:energy_intensity}





\cleardoublepage
           % \ref{chap:energy_intensity}
    
% The next command tells RStudio to do "Compile PDF" on book.Rnw,
% instead of this chapter, thereby eliminating the need to switch back to book.Rnw 
% before making the book.
%!TEX root = ../../book.Rnw

\chapter{Carbon intensity of energy}
\label{chap:carbon_intensity}





\cleardoublepage
           % \ref{chap:chap:carbon_intensity}
    
% The next command tells RStudio to do "Compile PDF" on book.Rnw,
% instead of this chapter, thereby eliminating the need to switch back to book.Rnw 
% before making the book.
%!TEX root = ../../book.Rnw

\chapter{Households}
\label{chap:households}





\cleardoublepage
                     % \ref{chap:households}
    
% The next command tells RStudio to do "Compile PDF" on book.Rnw,
% instead of this chapter, thereby eliminating the need to switch back to book.Rnw 
% before making the book.
%!TEX root = ../../book.Rnw

\chapter{Transportation}
\label{chap:transportation}





\cleardoublepage
             % \ref{chap:transportation}
    
% The next command tells RStudio to do "Compile PDF" on book.Rnw,
% instead of this chapter, thereby eliminating the need to switch back to book.Rnw 
% before making the book.
%!TEX root = ../../book.Rnw

\chapter{Agriculture}
\label{chap:agriculture}





\cleardoublepage
                   % \ref{chap:agriculture}
    
% The next command tells RStudio to do "Compile PDF" on book.Rnw,
% instead of this chapter, thereby eliminating the need to switch back to book.Rnw 
% before making the book.
%!TEX root = ../../book.Rnw

\chapter{Land use and urban planning}
\label{chap:land_use_urban_planning}





\cleardoublepage
 % \ref{chap:land_use_urban_planning}
    
% The next command tells RStudio to do "Compile PDF" on book.Rnw,
% instead of this chapter, thereby eliminating the need to switch back to book.Rnw 
% before making the book.
%!TEX root = ../../book.Rnw

\chapter{Government}
\label{chap:government}





\cleardoublepage
                     % \ref{chap:government}
    
% The next command tells RStudio to do "Compile PDF" on book.Rnw,
% instead of this chapter, thereby eliminating the need to switch back to book.Rnw 
% before making the book.
%!TEX root = ../../book.Rnw

\chapter{Personal action}
\label{chap:personal_action}





\cleardoublepage
             % \ref{chap:personal_action}

 %   \input{biblio}              % bibliography

    \bibliographystyle{unsrtnat}
    \bibliography{MCBook2021}   %chktex 11
    
    
    \cleardoublepage

 \backmatter                    % back matter
    \input{bio}                 % author(s)'s short bio(s)
								% please see the bios folder for some pdf examples
    
    \cleardoublepage

    \input{book.ind}            % if you have an Index

\end{document}

            % if you have an Index

\end{document}

            % if you have an Index

\end{document}

            % if you have an Index

\end{document}

